% Generated by the blueprint tool (blueprint-pdf/1) from a blueprint snapshot; edit the
% Markdown sources, not this file.  Compile with `latexmk -pdf` or by running
% `pdflatex` twice (the second run settles the contents and the page references).
% Images, if any, are read from assets/ next to this file.
\documentclass[11pt]{article}
\usepackage[T1]{fontenc}
\usepackage[utf8]{inputenc}
\usepackage{lmodern}
\usepackage{textcomp}
\usepackage[a4paper,margin=2.5cm,headheight=14pt]{geometry}
\usepackage{amsmath,amssymb,amsthm}
\usepackage{amscd}  % KaTeX's CD diagrams; part of amsmath, so always there
\usepackage{mathtools}
\usepackage{bm}
\usepackage{graphicx}
\usepackage{xcolor}
\usepackage{array,booktabs,tabularx}
\usepackage{fancyvrb}
\usepackage[normalem]{ulem}
\usepackage{microtype}
% Optional, and only with their fonts: a small TeX installation has the
% packages without the fonts, which fails at the first formula.
\IfFileExists{mathrsfs.sty}{\IfFileExists{rsfs10.tfm}{\usepackage{mathrsfs}}{}}{}
\IfFileExists{stmaryrd.sty}{\IfFileExists{stmary10.tfm}{\usepackage{stmaryrd}}{}}{}
\IfFileExists{cancel.sty}{\usepackage{cancel}}{}
\usepackage[hyphens]{url}
\usepackage[pdfusetitle,bookmarksnumbered=false]{hyperref}
\usepackage{bookmark}
\definecolor{bplink}{HTML}{1F4E9C}
\definecolor{bpmuted}{HTML}{5F676E}
\hypersetup{colorlinks=true,linkcolor=bplink,urlcolor=bplink,citecolor=bplink}
\setlength{\parindent}{0pt}
\setlength{\parskip}{0.55em plus 0.2em minus 0.1em}
% Lists as tight as the page's.
\makeatletter
\def\@listi{\leftmargin\leftmargini\parsep 0pt\topsep 0.3em plus 0.1em\itemsep 0.1em plus 0.1em}
\let\@listI\@listi
\def\@listii{\leftmargin\leftmarginii\labelwidth\leftmarginii\advance\labelwidth-\labelsep\parsep 0pt\topsep 0.1em\itemsep 0.1em}
% An ordered list that does not start at 1.
\newcommand{\bpstartat}[1]{\setcounter{\@enumctr}{\numexpr#1-1\relax}}
\makeatother
\setcounter{tocdepth}{3}
\emergencystretch=3em
\allowdisplaybreaks

% A numbered statement, set as a paper sets one: "Theorem 2.1 (Title)."
% #1 is the lead ("Theorem~2.1"), #2 what follows it in roman.
\newenvironment{bpstatement}[2]{\par\addvspace{0.6\baselineskip}\noindent\textbf{#1}#2\textbf{.}\enspace\ignorespaces}
  {\par\addvspace{0.4\baselineskip}}
\newenvironment{bpstatementit}[2]{\begin{bpstatement}{#1}{#2}\itshape\ignorespaces}{\end{bpstatement}}
% A step: a relation with prose of its own, read under the object it starts from.
\newenvironment{bpstep}[1]{\par\addvspace{0.3\baselineskip}\begin{list}{}{\leftmargin=1.6em\rightmargin=0pt\topsep=0pt}\item\relax
  {\small\color{bpmuted}#1}\par\nopagebreak}{\end{list}\par}
% Notes about where an object is written out.
\newenvironment{bpnote}{\par\small\color{bpmuted}\itshape}{\par}
\newcommand{\bpomitted}[1]{\par{\small\color{bpmuted}\itshape #1\par}}
\newcommand{\bpbodyheading}[1]{\par\addvspace{0.4\baselineskip}\noindent\textbf{#1}\par\nopagebreak}
\newcommand{\bprule}{\par\medskip\noindent\hfil\rule{0.3\linewidth}{0.4pt}\hfil\par\medskip}
\newcommand{\bpboundary}[1]{\par\nopagebreak{\small\color{bpmuted}#1\par}\nopagebreak}
% A reference to something this export leaves out, and one that resolves nowhere.
\newcommand{\bpout}[1]{#1\textsuperscript{\hyperlink{bp-outside}{\ensuremath{\ast}}}}
\newcommand{\bpbroken}[1]{\textcolor{red!70!black}{[#1?]}}
\newcommand{\bpslug}[1]{\texttt{#1}}
% A bibliography entry: label, anchor, entry; the label hangs in the margin.
\newcommand{\bpbibitem}[3]{\par\setbox0\hbox{[#1]\enspace}\ifdim\wd0<5em\setbox0\hbox to 5em{[#1]\hfil}\fi
  \hangindent=\wd0 \hangafter=1 \noindent\box0\phantomsection\label{#2}#3\par\smallskip}
% A character the export knows no way to set.
\newcommand{\bpglyph}[1]{\fbox{\scriptsize #1}}
% Images, at their own size unless they are wider than the line.
\newsavebox{\bpimagebox}
\newcommand{\bpimage}[1]{\sbox{\bpimagebox}{\includegraphics{#1}}%
  \ifdim\wd\bpimagebox>\linewidth\resizebox{\linewidth}{!}{\usebox{\bpimagebox}}\else\usebox{\bpimagebox}\fi}
\newcommand{\bpmissingimage}[1]{\fbox{\parbox{0.8\linewidth}{\centering\small\color{bpmuted}Image not included: #1}}}
\DefineVerbatimEnvironment{bpverbatim}{Verbatim}{fontsize=\small,xleftmargin=1em}
\renewcommand{\qedsymbol}{\ensuremath{\blacksquare}}
\hypersetup{pdftitle={Weil Conjectures — 8 Smooth base change},pdfsubject={Chapter 8 Smooth base change; all levels},pdfcreator={blueprint}}
% KaTeX commands the website knows
\providecommand{\R}{\mathbb{R}}\providecommand{\N}{\mathbb{N}}\providecommand{\Z}{\mathbb{Z}}
\providecommand{\Q}{\mathbb{Q}}\providecommand{\C}{\mathbb{C}}
\providecommand{\Reals}{\mathbb{R}}\providecommand{\reals}{\mathbb{R}}\providecommand{\Complex}{\mathbb{C}}
\providecommand{\cnums}{\mathbb{C}}\providecommand{\natnums}{\mathbb{N}}\providecommand{\Bbb}[1]{\mathbb{#1}}
\providecommand{\bold}[1]{\mathbf{#1}}\providecommand{\frak}[1]{\mathfrak{#1}}\providecommand{\mathbbm}[1]{\mathbb{#1}}
\providecommand{\lang}{\langle}\providecommand{\rang}{\rangle}\providecommand{\bra}[1]{\langle #1\rvert}
\providecommand{\ket}[1]{\lvert #1\rangle}\providecommand{\braket}[1]{\langle #1\rangle}
\providecommand{\set}[1]{\{#1\}}\providecommand{\Set}[1]{\left\{#1\right\}}
\providecommand{\infin}{\infty}\providecommand{\isin}{\in}\providecommand{\plusmn}{\pm}\providecommand{\sub}{\subset}
\providecommand{\sube}{\subseteq}\providecommand{\supe}{\supseteq}\providecommand{\exist}{\exists}
\providecommand{\bull}{\bullet}\providecommand{\larr}{\leftarrow}\providecommand{\rarr}{\rightarrow}
\providecommand{\lrarr}{\leftrightarrow}\providecommand{\Larr}{\Leftarrow}\providecommand{\Rarr}{\Rightarrow}
\providecommand{\Lrarr}{\Leftrightarrow}\providecommand{\uarr}{\uparrow}\providecommand{\darr}{\downarrow}
\providecommand{\harr}{\leftrightarrow}\providecommand{\lArr}{\Leftarrow}\providecommand{\rArr}{\Rightarrow}
\providecommand{\hArr}{\Leftrightarrow}\providecommand{\alef}{\aleph}\providecommand{\alefsym}{\aleph}
\providecommand{\real}{\Re}\providecommand{\image}{\Im}\providecommand{\weierp}{\wp}
\providecommand{\Alpha}{\mathrm{A}}\providecommand{\Beta}{\mathrm{B}}\providecommand{\Epsilon}{\mathrm{E}}
\providecommand{\Zeta}{\mathrm{Z}}\providecommand{\Eta}{\mathrm{H}}\providecommand{\Iota}{\mathrm{I}}
\providecommand{\Kappa}{\mathrm{K}}\providecommand{\Mu}{\mathrm{M}}\providecommand{\Nu}{\mathrm{N}}
\providecommand{\Omicron}{\mathrm{O}}\providecommand{\Rho}{\mathrm{P}}\providecommand{\Tau}{\mathrm{T}}
\providecommand{\Chi}{\mathrm{X}}\providecommand{\omicron}{o}\providecommand{\gt}{>}\providecommand{\lt}{<}
\providecommand{\And}{\&}\providecommand{\hdashline}{\hline}\providecommand{\operatornamewithlimits}[1]{\operatorname*{#1}}
\providecommand{\llbracket}{[\![}\providecommand{\rrbracket}{]\!]}\providecommand{\cancel}[1]{#1}
\providecommand{\mathscr}{\mathcal}
% the characters beyond ASCII this document uses
\DeclareUnicodeCharacter{2160}{\textup{I}}
\DeclareUnicodeCharacter{2161}{\textup{II}}
% the project's macros ([katex.macros] in blueprint.toml)
\providecommand{\A}{}\renewcommand{\A}{\mathbf A}
\providecommand{\AEn}{}\renewcommand{\AEn}{A_{E,n}}
\providecommand{\Aaf}{}\renewcommand{\Aaf}{\mathbf A}
\providecommand{\Ab}{}\renewcommand{\Ab}{\mathsf{Ab}}
\providecommand{\AffSch}{}\renewcommand{\AffSch}{\operatorname{AffSch}}
\providecommand{\Aone}{}\renewcommand{\Aone}{\mathbf A^1}
\providecommand{\Bl}{}\renewcommand{\Bl}{\operatorname{Bl}}
\providecommand{\Br}{}\renewcommand{\Br}{\mathrm{Br}}
\providecommand{\C}{}\renewcommand{\C}{\mathbb C}
\providecommand{\Cech}{}\renewcommand{\Cech}{\check{H}}
\providecommand{\Coext}{}\renewcommand{\Coext}{\operatorname{Coext}}
\providecommand{\Cone}{}\renewcommand{\Cone}{\operatorname{Cone}}
\providecommand{\Cons}{}\renewcommand{\Cons}{\operatorname{Cons}}
\providecommand{\Cref}{}\renewcommand{\Cref}[1]{\text{#1}}
\providecommand{\D}{}\renewcommand{\D}{D}
\providecommand{\DLAFT}{}\renewcommand{\DLAFT}{\operatorname{FT}}
\providecommand{\DLAFr}{}\renewcommand{\DLAFr}{\operatorname{Fr}}
\providecommand{\DLAQ}{}\renewcommand{\DLAQ}{{\overline{\mathbf Q}_\ell}}
\providecommand{\DLRQ}{}\renewcommand{\DLRQ}{{\overline{\mathbf Q}_\ell}}
\providecommand{\Dbc}{}\renewcommand{\Dbc}{D^b_c}
\providecommand{\Dc}{}\renewcommand{\Dc}{D_{\mathrm{cons}}}
\providecommand{\Dctf}{}\renewcommand{\Dctf}{D^b_{ctf}}
\providecommand{\Dic}{}\renewcommand{\Dic}{D_{\IndCons}}
\providecommand{\Dplus}{}\renewcommand{\Dplus}{D^{+}}
\providecommand{\End}{}\renewcommand{\End}{\operatorname{End}}
\providecommand{\Et}{}\renewcommand{\Et}{\mathrm{\acute{e}t}}
\providecommand{\Ex}{}\renewcommand{\Ex}{\mathrm{Ex}}
\providecommand{\Ext}{}\renewcommand{\Ext}{\operatorname{Ext}}
\providecommand{\F}{}\renewcommand{\F}{\mathbb F}
\providecommand{\FT}{}\renewcommand{\FT}{\operatorname{FT}}
\providecommand{\Fbar}{}\renewcommand{\Fbar}{{\overline{\mathbf F}_q}}
\providecommand{\Fell}{}\renewcommand{\Fell}{\mathbf F_\ell}
\providecommand{\Fib}{}\renewcommand{\Fib}{\operatorname{Fib}}
\providecommand{\Fix}{}\renewcommand{\Fix}{\operatorname{Fix}}
\providecommand{\Fp}{}\renewcommand{\Fp}{{\mathbf F_{p}}}
\providecommand{\Fq}{}\renewcommand{\Fq}{{\mathbf F_q}}
\providecommand{\Fqn}{}\renewcommand{\Fqn}{{\mathbf F_{q^{n}}}}
\providecommand{\Frob}{}\renewcommand{\Frob}{\operatorname{Frob}}
\providecommand{\GL}{}\renewcommand{\GL}{\operatorname{GL}}
\providecommand{\Gal}{}\renewcommand{\Gal}{\operatorname{Gal}}
\providecommand{\Hens}{}\renewcommand{\Hens}{\operatorname{Hens}}
\providecommand{\Hilb}{}\renewcommand{\Hilb}{\operatorname{Hilb}}
\providecommand{\Hom}{}\renewcommand{\Hom}{\operatorname{Hom}}
\providecommand{\IndCons}{}\renewcommand{\IndCons}{\operatorname{IndCons}}
\providecommand{\Kar}{}\renewcommand{\Kar}{\operatorname{Kar}}
\providecommand{\LT}{}\renewcommand{\LT}{\operatorname{LT}}
\providecommand{\Lie}{}\renewcommand{\Lie}{\operatorname{Lie}}
\providecommand{\Loc}{}\renewcommand{\Loc}{\operatorname{Loc}}
\providecommand{\Mod}{}\renewcommand{\Mod}{\mathsf{Mod}}
\providecommand{\N}{}\renewcommand{\N}{\mathbb{N}}
\providecommand{\Needspace}{}\renewcommand{\Needspace}{\par}
\providecommand{\O}{}\renewcommand{\O}{{\mathcal{O}}}
\providecommand{\Ob}{}\renewcommand{\Ob}{\mathrm{Ob}}
\providecommand{\Oe}{}\renewcommand{\Oe}{\mathcal O_E}
\providecommand{\PP}{}\renewcommand{\PP}{\mathbf{P}}
\providecommand{\Perf}{}\renewcommand{\Perf}{\operatorname{Perf}}
\providecommand{\Pic}{}\renewcommand{\Pic}{\mathrm{Pic}}
\providecommand{\Pm}{}\renewcommand{\Pm}{\mathbf P^{m}}
\providecommand{\Pone}{}\renewcommand{\Pone}{\mathbf P^{1}}
\providecommand{\Pqn}{}\renewcommand{\Pqn}{\mathbf P^{n}}
\providecommand{\ProFin}{}\renewcommand{\ProFin}{\operatorname{ProFin}}
\providecommand{\Proet}{}\renewcommand{\Proet}{\mathrm{pro\acute{e}t}}
\providecommand{\Q}{}\renewcommand{\Q}{\mathbb Q}
\providecommand{\Qb}{}\renewcommand{\Qb}{\overline{\mathbf Q}_{\ell}}
\providecommand{\Qbar}{}\renewcommand{\Qbar}{\overline{\mathbf Q}}
\providecommand{\QbarX}{}\renewcommand{\QbarX}{\overline{\mathbf Q}_{\ell,X}}
\providecommand{\Qbarl}{}\renewcommand{\Qbarl}{{\overline{\mathbf Q}_{\ell}}}
\providecommand{\Qell}{}\renewcommand{\Qell}{{\mathbf{Q}_{\ell}}}
\providecommand{\Ql}{}\renewcommand{\Ql}{\mathbf Q_{\ell}}
\providecommand{\Qlb}{}\renewcommand{\Qlb}{{\overline{\mathbf Q}_\ell}}
\providecommand{\Qlbar}{}\renewcommand{\Qlbar}{\overline{\mathbf Q}_\ell}
\providecommand{\R}{}\renewcommand{\R}{\mathbb{R}}
\providecommand{\RG}{}\renewcommand{\RG}{\operatorname{\mathrm{R}\Gamma}}
\providecommand{\RGamma}{}\renewcommand{\RGamma}{R\Gamma}
\providecommand{\RGammac}{}\renewcommand{\RGammac}{R\Gamma_c}
\providecommand{\RHom}{}\renewcommand{\RHom}{R\mathcal{H}om}
\providecommand{\RHomX}{}\renewcommand{\RHomX}{R\mathcal H om_X}
\providecommand{\Rat}{}\renewcommand{\Rat}{\operatorname{IsRat}}
\providecommand{\Res}{}\renewcommand{\Res}{\operatorname{Res}}
\providecommand{\Rlim}{}\renewcommand{\Rlim}{R\!\varprojlim}
\providecommand{\SIm}{}\renewcommand{\SIm}{\operatorname{SIm}}
\providecommand{\Sch}{}\renewcommand{\Sch}{\mathsf{Sch}}
\providecommand{\Sh}{}\renewcommand{\Sh}{\operatorname{Sh}}
\providecommand{\Span}{}\renewcommand{\Span}{\operatorname{span}}
\providecommand{\Spec}{}\renewcommand{\Spec}{\operatorname{Spec}}
\providecommand{\Tr}{}\renewcommand{\Tr}{\operatorname{Tr}}
\providecommand{\Trsm}{}\renewcommand{\Trsm}{\operatorname{Tr}^{\mathrm{sm}}}
\providecommand{\Wt}{}\renewcommand{\Wt}{\operatorname{Wt}}
\providecommand{\Z}{}\renewcommand{\Z}{\mathbb{Z}}
\providecommand{\Zar}{}\renewcommand{\Zar}{\mathrm{zar}}
\providecommand{\Zell}{}\renewcommand{\Zell}{{\mathbf{Z}_{\ell}}}
\providecommand{\Zl}{}\renewcommand{\Zl}{{\mathbf Z_\ell}}
\providecommand{\affet}{}\renewcommand{\affet}[1]{{#1}^{\mathrm{aff}}_{\mathrm{ét}}}
\providecommand{\affproet}{}\renewcommand{\affproet}[1]{{#1}^{\mathrm{aff}}_{\mathrm{proét}}}
\providecommand{\boxt}{}\renewcommand{\boxt}{\boxtimes_A^{\mathrm L}}
\providecommand{\cF}{}\renewcommand{\cF}{\mathcal{F}}
\providecommand{\cHom}{}\renewcommand{\cHom}{R\mathcal{H}\!om}
\providecommand{\cI}{}\renewcommand{\cI}{\mathcal{I}}
\providecommand{\cL}{}\renewcommand{\cL}{\mathcal L}
\providecommand{\cM}{}\renewcommand{\cM}{\mathcal M}
\providecommand{\cN}{}\renewcommand{\cN}{\mathcal N}
\providecommand{\cO}{}\renewcommand{\cO}{\mathcal O}
\providecommand{\cT}{}\renewcommand{\cT}{\mathcal T}
\providecommand{\checked}{}\renewcommand{\checked}{\color{red}\aftergroup\restorecolor}
\providecommand{\coeff}{}\renewcommand{\coeff}{\operatorname{coeff}}
\providecommand{\coker}{}\renewcommand{\coker}{\operatorname{coker}}
\providecommand{\colim}{}\renewcommand{\colim}{\operatorname{colim}}
\providecommand{\constructedby}{}\renewcommand{\constructedby}[1]{\par\smallskip\noindent\textcolor{purple!60!black}{\textsf{Constructed by:} }
  \footnotesize\sloppy\nolinkurl{#1}\par}
\providecommand{\deepresult}{}\renewcommand{\deepresult}[1]{\par\smallskip\noindent\textcolor{orange!45!black}{
  \textsf{Blueprint role: named deep theorem node; see the cluster proof architecture.}}
  \par\noindent\textcolor{orange!45!black}{\textsf{Mathematical source:} }#1\par}
\providecommand{\dt}{}\renewcommand{\dt}{\mathbin{\otimes^{\mathbf L}}}
\providecommand{\dw}{}\renewcommand{\dw}{\operatorname{dw}}
\providecommand{\et}{}\renewcommand{\et}[1]{{#1}_{\mathrm{ét}}}
\providecommand{\ev}{}\renewcommand{\ev}{\operatorname{ev}}
\providecommand{\fib}{}\renewcommand{\fib}{\operatorname{fib}}
\providecommand{\formalizationnote}{}\renewcommand{\formalizationnote}[1]{\begin{quote}\small\raggedright\textbf{Formalization note.} #1\end{quote}}
\providecommand{\fpqc}{}\renewcommand{\fpqc}{X_{\mathrm{fpqc}}}
\providecommand{\gammacont}{}\renewcommand{\gammacont}{\Gamma_{\mathrm{cont}}}
\providecommand{\gm}{}\renewcommand{\gm}{\mathbf G_m}
\providecommand{\hocolim}{}\renewcommand{\hocolim}{\operatorname{hocolim}}
\providecommand{\holim}{}\renewcommand{\holim}{\operatorname{holim}}
\providecommand{\hrB}{}\renewcommand{\hrB}{\@ifnextchar]{\hrEnd}{\hrStep}}
\providecommand{\hrEnd}{}\renewcommand{\hrEnd}{\@firstoftwo{}}
\providecommand{\hrStep}{}\renewcommand{\hrStep}{\@firstoftwo{\hrB}}
\providecommand{\id}{}\renewcommand{\id}{\operatorname{id}}
\providecommand{\kunFl}{}\renewcommand{\kunFl}{\mathbf F_\ell}
\providecommand{\kunZn}{}\renewcommand{\kunZn}{\mathbf Z/n\mathbf Z}
\providecommand{\kuncd}{}\renewcommand{\kuncd}{\operatorname{cd}}
\providecommand{\kundim}{}\renewcommand{\kundim}{\dim}
\providecommand{\kunotimes}{}\renewcommand{\kunotimes}[1]{\otimes^{\mathbf L}_{#1}}
\providecommand{\kunsupp}{}\renewcommand{\kunsupp}{\operatorname{Supp}}
\providecommand{\lotimes}{}\renewcommand{\lotimes}{\otimes^{\mathrm L}}
\providecommand{\m}{}\renewcommand{\m}{{\mathfrak{m}}}
\providecommand{\n}{}\renewcommand{\n}{{\mathfrak{n}}}
\providecommand{\nagataGap}{}\renewcommand{\nagataGap}[1]{\par\noindent\textbf{Formalization gap.} #1\par}
\providecommand{\openobligation}{}\renewcommand{\openobligation}[1]{\par\smallskip\noindent\textcolor{red!65!black}{\textsf{Open obligation:} }
  \footnotesize\sloppy\nolinkurl{#1}\par}
\providecommand{\p}{}\renewcommand{\p}{{\mathfrak{p}}}
\providecommand{\pr}{}\renewcommand{\pr}{\operatorname{pr}}
\providecommand{\preshv}{}\renewcommand{\preshv}[1]{\mathrm{PreShv}(#1)}
\providecommand{\pro}{}\renewcommand{\pro}[1]{\mathrm{Pro}(#1)}
\providecommand{\proet}{}\renewcommand{\proet}{X_{\mathrm{proét}}}
\providecommand{\q}{}\renewcommand{\q}{{\mathfrak{q}}}
\providecommand{\rad}{}\renewcommand{\rad}{\operatorname{rad}}
\providecommand{\rank}{}\renewcommand{\rank}{\operatorname{rank}}
\providecommand{\red}{}\renewcommand{\red}{\color{red}}
\providecommand{\restorecolor}{}\renewcommand{\restorecolor}{\color{black}}
\providecommand{\retiredobligation}{}\renewcommand{\retiredobligation}[1]{\par\smallskip\noindent\textcolor{red!75!black}{\textsf{Retired invalid route:} }
  \footnotesize\sloppy\nolinkurl{#1}\par}
\providecommand{\rk}{}\renewcommand{\rk}{\operatorname{rk}}
\providecommand{\sep}{}\renewcommand{\sep}{\textnormal{sep}}
\providecommand{\shv}{}\renewcommand{\shv}[1]{\mathrm{Shv}(#1)}
\providecommand{\spec}{}\renewcommand{\spec}[1]{\mathrm{Spec}(#1)}
\providecommand{\stacks}{}\renewcommand{\stacks}[1]{\href{https://stacks.math.columbia.edu/tag/#1}{\textnormal{Stacks, Tag~#1}}}
\providecommand{\stacksproject}{}\renewcommand{\stacksproject}[1]{{\cite[\href{https://stacks.math.columbia.edu/tag/#1}{Tag #1}]{stacks-project}}}
\providecommand{\statusnote}{}\renewcommand{\statusnote}[1]{\begin{quote}\small\raggedright\textbf{Status.} #1\end{quote}}
\providecommand{\tildeideal}{}\renewcommand{\tildeideal}[2]{\widetilde{(#1, #2)}}
\providecommand{\tildering}{}\renewcommand{\tildering}[3]{{#1}_{\widetilde{(#2, #3)}}}
\providecommand{\tr}{}\renewcommand{\tr}{\operatorname{tr}}
\providecommand{\wt}{}\renewcommand{\wt}{\operatorname{wt}}

\begin{document}
\pagestyle{plain}
\begin{center}
{\LARGE\bfseries Weil Conjectures\par}
\vspace{0.7em}
{\large Chapter 8: Smooth base change\par}
\vspace{0.5em}
{\small\color{bpmuted}Snapshot 2026-09-26, revision 8e4cb88e2a\par}
\end{center}
\vspace{0.5em}
{\small\tableofcontents}
\bigskip\hrule\bigskip

% object: chap-Smooth-base-change (blueprint/chap-Smooth-base-change/_section.md)
\section*{\texorpdfstring{8\quad Smooth base change}{8 Smooth base change}}
\phantomsection\label{bp1}%
\addcontentsline{toc}{section}{\texorpdfstring{8\quad Smooth base change}{8 Smooth base change}}
% text: prose-08-01
\phantomsection\label{bp2}%
To construct the trace classes in the trace formula we need the Künneth formula for ordinary direct images, \(Rf_{1*}M_1\boxtimes Rf_{2*}M_2\simeq R(f_1\times f_2)_*(M_1\boxtimes M_2)\) (Proposition \hyperref[bp64]{8.18}, Section \hyperref[bp49]{8.9}). For \(Rf_!\) this formula follows from proper base change and the projection formula, but for \(Rf_*\) this argument does not work, because \(Rf_*\) does not commute with base change in general. For example, let \(j:U=\mathbf A^1\setminus\{0\}\hookrightarrow\mathbf A^1\) over an algebraically closed field, and let \(n\) be invertible in it. Then the stalk of \(R^1j_*(\mathbf Z/n)\) at \(0\) is \(H^1\) of the punctured strictly henselian disc, which is \(\mathbf Z/n(-1)\neq0\), whereas the fibre of \(j\) over \(0\) is empty, so base change along the closed immersion of \(0\) fails.

The Künneth formula only needs base change along a product projection \(Y\times T\to Y\) (Lemma \hyperref[bp51]{8.13}). The question is whether base change along such a map holds for every \(f\).

In topology it does, at least when \(T\) is a manifold. If \(g\) is a local homeomorphism, base change along \(g\) is restriction to open subsets. If \(g\) is the projection \(Y\times\mathbf R^m\to Y\), then the stalk of \(R^qf'_*\) at \((y,t)\) is the colimit of the cohomology groups of \(f^{-1}(V)\times B\) over small neighbourhoods \(V\times B\) of \((y,t)\), and the contractible ball \(B\) does not change them. For schemes the analogue is less clear: a smooth morphism of schemes is, Zariski locally on the source, étale over some \(\mathbf A^d_Y\), but the affine line has no small contractible neighbourhoods. Nevertheless, the theorem of this chapter, due to Artin (SGA 4, Exposé XVI [\hyperref[bp66]{SGA4}]), says that for coefficients prime to the characteristic the affine line behaves as if it had them. More precisely, let \(g:Y'\to Y\) be smooth, let \(f\) be quasi-compact and quasi-separated, and let \(F\) be a torsion sheaf whose torsion orders are invertible on \(Y\). Then the base change map \(g^*R^qf_*F\to R^qf'_*g'^*F\) is an isomorphism (Theorem \hyperref[bp7]{8.1}). Moreover, on \(D^b_{\mathrm{ctf}}(-,\Lambda)\) the same holds for every admissible \(\Lambda\) (Theorem \hyperref[bp45]{8.11}).

A product projection \(Y\times T\to Y\) is not smooth when \(T\) is singular. What the Künneth formula uses is a consequence of smooth base change, namely base change under extension of the ground field [Stacks, \href{https://stacks.math.columbia.edu/tag/0F1C}{0F1C}].

The essential case is a change of the algebraically closed ground field, and there the key input is a fact about curves. Let \(\Omega/L\) be an extension of algebraically closed fields, \(C\) a smooth affine curve over \(L\), and \(n\) invertible in \(L\). Then pullback \(H^1(C,\mathbf Z/n)\to H^1(C_\Omega,\mathbf Z/n)\) is an isomorphism (Theorem \bpout{4.20}). In other words, a cyclic cover of \(C_\Omega\) of degree \(n\) does not depend on the extra transcendental parameters of \(\Omega\).

For \(p\)-torsion coefficients this invariance fails, and so does smooth base change. Let \(k\) be algebraically closed of characteristic \(p>0\), and let \(f:\mathbf A^1_k\to\operatorname{Spec}k\) and \(g:\mathbf A^1_k\to\operatorname{Spec}k\) be the structure maps, with coordinates \(x\) on the source of \(f\) and \(t\) on the source of \(g\). By the Artin–Schreier sequence, \(H^1(\mathbf A^1_K,\mathbf F_p)\) is the cokernel of \(h\mapsto h^p-h\) on \(K[x]\), for every field \(K\supset k\). At the geometric generic point of the \(t\)-line, the stalk of \(g^*R^1f_*\mathbf F_p\) is \(H^1(\mathbf A^1_k,\mathbf F_p)\), the stalk of \(R^1f'_*\mathbf F_p\) is \(H^1(\mathbf A^1_K,\mathbf F_p)\) with \(K=k(t)^{\mathrm{sep}}\), and the base change map is pullback. We claim that the class of \(tx\) is not in its image. Suppose \(tx=a+h^p-h\) with \(a\in k[x]\) and \(h\in K[x]\), and let \(b_e\) be the coefficient of \(x^{p^e}\) in \(h\). Comparing the coefficients of \(x\) gives \(b_0\in -t+k\), and comparing those of \(x^{p^{e+1}}\) gives \(b_{e+1}\in b_e^p+k\). By induction \(b_e\in -t^{p^e}+k\), so \(b_e\neq0\) for every \(e\geq0\). This is absurd for a polynomial. In geometric terms, the Artin–Schreier cover \(y^p-y=tx\) depends on the parameter \(t\). Hence the hypothesis on the torsion orders cannot be dropped.

% object: sec-The-base-change-map-2 (blueprint/chap-Smooth-base-change/sec-The-base-change-map-2/_section.md), content hidden
\subsection*{\texorpdfstring{8.1\quad The base change map}{8.1 The base change map}}
\phantomsection\label{bp3}%
\addcontentsline{toc}{subsection}{\texorpdfstring{8.1\quad The base change map}{8.1 The base change map}}
% text: prose-08-02
\phantomsection\label{bp4}%
Consider a cartesian square, with \(X'=Y'\times_YX\),

\[
\begin{CD}
X' @>{g'}>> X\\
@V{f'}VV @VV{f}V\\
Y' @>{g}>> Y.
\end{CD}
\tag{8.1}
\]

The base change map \(g^*Rf_*K\to Rf'_*g'^*K\), for \(K\in D(X_{\mathrm{pro}\acute{e}t},\Lambda)\), is the map (4.5) of Section \bpout{4.1}. It comes by adjunction from the counit of \(f^*\dashv Rf_*\), and it is defined for every such square and every complex, with \(g\) smooth or not. If \(f\) is quasi-compact and quasi-separated (qcqs), its stalk at a geometric point \(\bar y'\) of \(Y'\) with image \(\bar y\) in \(Y\) is the pullback map from the cohomology of \(X\times_Y\operatorname{Spec}\mathcal O^{\mathrm{sh}}_{Y,\bar y}\) to that of \(X'\times_{Y'}\operatorname{Spec}\mathcal O^{\mathrm{sh}}_{Y',\bar y'}\).

In smooth base change for torsion sheaves (Theorem \hyperref[bp7]{8.1}) we assume that \(f\) is qcqs and that \(g\) is smooth. In the theorem on \(D^b_{\mathrm{ctf}}\) all schemes are of finite type over \(k\). Every inverse image below is the exact inverse image for the pullback of the coefficient ring, so that \(g^*=Lg^*\) and \(g'^*=Lg'^*\).

The proofs of these theorems transport the base change map in three ways: from the pro-étale to the étale site, where the classical arguments are made; along composites of squares, when a smooth morphism is factored into curves; and along a change of the coefficient ring, to reach \(\ell\)-adic coefficients. Since the map is built from units and counits, it is compatible with all three, by Lemma \bpout{4.7} for the first two and by Lemma \hyperref[bp43]{8.10} in Section \hyperref[bp41]{8.8} below for the third.


% object: sec-Torsion-sheaves (blueprint/chap-Smooth-base-change/sec-Torsion-sheaves/_section.md), content hidden
\subsection*{\texorpdfstring{8.2\quad Torsion sheaves}{8.2 Torsion sheaves}}
\phantomsection\label{bp5}%
\addcontentsline{toc}{subsection}{\texorpdfstring{8.2\quad Torsion sheaves}{8.2 Torsion sheaves}}
% text: prose-08-03
\phantomsection\label{bp6}%
Smooth base change for complexes (Section \hyperref[bp36]{8.7}) and for \(\ell\)-adic coefficients (Section \hyperref[bp41]{8.8}) follows from the case of a single torsion sheaf, which is the next theorem. Its proof leaves schemes of finite type over \(k\): it passes through base schemes such as spectra of fields and of integral closures. The comparison results of Chapter \bpout{2} hold over arbitrary schemes for bounded below complexes, so they apply to these bases as well.

% object: theorem-smooth-base-change-classical-torsion (blueprint/chap-Smooth-base-change/sec-Torsion-sheaves/theorem-smooth-base-change-classical-torsion/_section.md)
\phantomsection\label{bp7}%
\begin{bpstatementit}{Theorem~8.1}{ (Smooth base change for classical torsion sheaves)}
In (8.1), suppose \(g\) is smooth and \(f\) is quasi-compact and quasi-separated. Let \(F\) be a classical torsion abelian sheaf on \(X_{\mathrm{pro}\acute{e}t}\) which is the filtered union of its subsheaves \(F[n]\) for the positive integers \(n\) invertible on \(Y\). Then for every \(q\geq0\) the degree-\(q\) base change map induced by (4.5)

\[
g^*R^qf_*F\longrightarrow R^qf'_*g'^*F
\]

is an isomorphism.
\end{bpstatementit}

% text: prose-08-04
\phantomsection\label{bp8}%
As for proper base change, we transport the problem to the étale site by the comparison theorems of Chapter \bpout{2}. Since \(F\) is classical, \(F=\nu_X^*F_{\acute{e}t}\) for a torsion abelian sheaf \(F_{\acute{e}t}\) on \(X_{\acute{e}t}\). As \(f\) and \(f'\) are qcqs, Theorem \bpout{2.38} and Lemma \bpout{4.7} identify both sides and the base change map with the inverse images of the corresponding étale objects. We may therefore work on the étale sites in the proof below, where \(g^{-1}\) denotes inverse image of étale sheaves.

We follow the second proof of [Stacks, \href{https://stacks.math.columbia.edu/tag/0EYU}{0EYU}], an induction on the cohomological degree \(q\). The idea is to concentrate a first failure of base change in a single sheaf \(R^qh_*\mathbf Z/m\), and then to show that this sheaf vanishes. In degree zero, by contrast, the base change map can be computed directly on sections. Moreover, base change may be checked étale locally and is stable under composition, so that étale coordinates reduce us to the case where \(g\) is a smooth affine curve.

Now suppose that base change holds in degrees below \(q\) but fails in degree \(q\). The failure persists when we shrink the source of the direct image by limit arguments, and also when we make integral changes of the source and of the base. In this way it is concentrated in a single situation, in which the coefficients are constant, equal to \(\mathbf Z/m\), and we have a smooth affine curve \(C\) over the spectrum \(S\) of a normal domain whose fraction field \(L\) is algebraically closed, together with the morphism \(h:C_\Omega\to C\) for an algebraically closed extension \(\Omega\) of \(L\). But in this situation \(R^qh_*\mathbf Z/m=0\). For \(q\geq2\) the reason is that affine curves over an algebraically closed field have no cohomology in degree at least two (Appendix \bpout{D}). For \(q=1\) it follows from the invariance of \(H^1\) of curves and an extension theorem for covers (Section \hyperref[bp30]{8.6}).


% object: sec-Degree-zero (blueprint/chap-Smooth-base-change/sec-Degree-zero/_section.md), content hidden
\subsection*{\texorpdfstring{8.3\quad Degree zero}{8.3 Degree zero}}
\phantomsection\label{bp9}%
\addcontentsline{toc}{subsection}{\texorpdfstring{8.3\quad Degree zero}{8.3 Degree zero}}
% text: prose-08-05
\phantomsection\label{bp10}%
In degree zero the base change map can be computed on sections, even for sheaves of sets. Sections of \(f_*F\) over \(V\) are sections of \(F\) over \(V\times_YX\), and the question is whether pulling back to a \(Y'\)-scheme \(U\) over \(V\) creates new sections. It does not when \(U\to V\) is flat, quasi-compact and quasi-separated with geometrically connected fibres, because pullback along such a map does not change sections of the pullback of an étale sheaf. Moreover, a flat \(g\) locally of finite presentation with geometrically reduced fibres factors locally into such a map followed by an étale one, and in this way Lemma \hyperref[bp11]{8.2} proves the isomorphism for every such \(g\).

% object: lemma-degree-zero-calculation (blueprint/chap-Smooth-base-change/sec-Degree-zero/lemma-degree-zero-calculation/_section.md)
\phantomsection\label{bp11}%
\begin{bpstatementit}{Lemma~8.2}{ (The degree-zero calculation)}
Suppose that \(g\) in (8.1) is flat and locally of finite presentation with geometrically reduced fibres, and that \(f\) is qcqs. Then the base change map \(g^{-1}f_*F_{\acute{e}t}\to f'_*g'^{-1}F_{\acute{e}t}\) is an isomorphism for every sheaf of sets \(F_{\acute{e}t}\) on \(X_{\acute{e}t}\).
\end{bpstatementit}

% proof: proof/lemma-degree-zero-calculation
\begin{proof}[{Proof}]\phantomsection\label{bp12}%
Both sides are sheaves on \(Y'_{\acute{e}t}\), so it is enough to compare their sections on a family of étale objects \(U\to Y'\) that cover \(Y'\). We choose these objects so that both sections can be computed from sections of \(F_{\acute{e}t}\) over a single scheme \(V\times_YX\).

Such objects exist locally on every étale \(U\to Y'\): by [Stacks, \href{https://stacks.math.columbia.edu/tag/0EY5}{0EY5}], locally on \(U\) there is a factorization \(U\to V\to Y\) in which \(V\to Y\) is étale and \(U\to V\) is flat and qcqs with geometrically connected fibres. Pullback along such a map does not change the sections of the pullback of an étale sheaf [Stacks, \href{https://stacks.math.columbia.edu/tag/0A3H}{0A3H}, \href{https://stacks.math.columbia.edu/tag/0EZX}{0EZX}], and we use this twice. Applying it to \(U\to V\), we compute the left side from \((f_*F_{\acute{e}t})(V)\); applying it to the base change \(U\times_YX\to V\times_YX\), which again has geometrically connected fibres, we compute the right side. On these objects we therefore have

\[
(g^{-1}f_*F_{\acute{e}t})(U)
=F_{\acute{e}t}(V\times_Y X)
=(f'_*g'^{-1}F_{\acute{e}t})(U),
\]

and the base change map is bijective on sections over every such \(U\), because these equalities are induced by it. As these \(U\) cover \(Y'\), the lemma follows by descent, as in the argument in degree zero of [Stacks, \href{https://stacks.math.columbia.edu/tag/0EYS}{0EYS}].
\end{proof}


% object: sec-The-induction-hypothesis (blueprint/chap-Smooth-base-change/sec-The-induction-hypothesis/_section.md), content hidden
\subsection*{\texorpdfstring{8.4\quad The induction hypothesis}{8.4 The induction hypothesis}}
\phantomsection\label{bp13}%
\addcontentsline{toc}{subsection}{\texorpdfstring{8.4\quad The induction hypothesis}{8.4 The induction hypothesis}}
% text: prose-08-06
\phantomsection\label{bp14}%
In positive degrees we argue by induction on the degree. This induction hypothesis has to survive the reductions outlined above: composition of base changes, étale localization, and replacement of the source and of the base of the direct image by other schemes. But a statement about the single square (8.1) does not survive them, so instead we fix an integer \(n\) invertible on \(Y\) and use a property of a morphism \(g:Y'\to Y\) which concerns all direct images over all \(Y\)-schemes at once, as in [Stacks, \href{https://stacks.math.columbia.edu/tag/0F02}{0F02}]. For \(r\geq0\), we say that \(\mathrm{BC}(g,n,r)\) holds if, for every scheme \(T\) over \(Y\), every qcqs morphism \(Z\to T\), every sheaf of \(\mathbf Z/n\)-modules \(M\) on \(Z_{\acute{e}t}\), and every \(0\leq i\leq r\), the base change map in degree \(i\) for the square with base change \(Y'\times_YT\to T\) and direct image along \(Z\to T\) is an isomorphism. We allow arbitrary schemes \(T\) because the reductions below lead to base schemes, such as the spectrum of a normal domain with algebraically closed fraction field, that are in general not Noetherian.

The property \(\mathrm{BC}(g,n,r)\) may be checked étale locally on \(Y\) and on \(Y'\). Moreover, for composable morphisms \(g_1\) and \(g_2\), \(\mathrm{BC}(g_1,n,r)\) and \(\mathrm{BC}(g_2,n,r)\) imply \(\mathrm{BC}(g_2\circ g_1,n,r)\). Both statements rest on two properties of the base change map in a fixed degree. First, the map is local for étale coverings of \(Y'\) and of \(Y\), since restriction to an étale object is localization of the site. Second, base change maps compose when two base changes are composed [Stacks, \href{https://stacks.math.columbia.edu/tag/0EZR}{0EZR}, \href{https://stacks.math.columbia.edu/tag/0EZS}{0EZS}].

For smooth \(g\), étale coordinates [Stacks, \href{https://stacks.math.columbia.edu/tag/054L}{054L}] give, locally on \(Y'\) and \(Y\) and with all schemes affine, an étale \(Y\)-morphism \(Y'\to\mathbf A^d_Y\), where \(d\) is the relative dimension of \(g\). Composing it with the projections \(\mathbf A^{j}_Y\to\mathbf A^{j-1}_Y\) that forget the last coordinate, we obtain a factorization

\[
Y'\longrightarrow\mathbf A^{d-1}_Y
\longrightarrow\cdots\longrightarrow\mathbf A^1_Y
\longrightarrow Y
\]

of \(g\) in which every displayed arrow is smooth of relative dimension one. Indeed, the first is an étale map followed by a projection, and the others are projections. The case \(d=0\) is étale localization. Hence, in the induction on the degree, it suffices to treat affine \(Y\) and \(Y'\) with \(g\) smooth of relative dimension one.

The coefficients reduce to \(\mathbf Z/n\)-modules. Indeed, a sheaf \(F\) as in Theorem \hyperref[bp7]{8.1} is the filtered union of the subsheaves \(F[n]\), ordered by divisibility of integers invertible on \(Y\), and inverse image commutes with this colimit, and so do higher direct images along qcqs morphisms [Stacks, \href{https://stacks.math.columbia.edu/tag/09Z1}{09Z1}]. It therefore suffices to treat \(\mathbf Z/n\)-modules for each such \(n\).


% object: sec-Reduction-to-curves (blueprint/chap-Smooth-base-change/sec-Reduction-to-curves/_section.md), content hidden
\subsection*{\texorpdfstring{8.5\quad Reduction to curves}{8.5 Reduction to curves}}
\phantomsection\label{bp15}%
\addcontentsline{toc}{subsection}{\texorpdfstring{8.5\quad Reduction to curves}{8.5 Reduction to curves}}
% text: prose-08-07
\phantomsection\label{bp16}%
Recall that \(\mathrm{BC}(g,n,r)\) asserts base change in degrees \(\leq r\), for \(\mathbf Z/n\)-coefficients, along every base change \(Y'\times_YT\to T\) of \(g\). Assume now that \(q\geq1\) and that \(\mathrm{BC}(g,n,q-1)\) holds. Under this assumption we show that if \(\mathrm{BC}(g,n,q)\) fails, it also fails in a special form: with constant coefficients and with direct image from the spectrum of an algebraically closed field (Lemma \hyperref[bp26]{8.6}). The next four lemmas, all on the étale sites, carry out this reduction step by step.

We argue by contradiction rather than by dévissage, because a nonzero cokernel of the base change map can be transported along an embedding of coefficient sheaves \(M\hookrightarrow N\), whereas an isomorphism cannot.

% object: lemma-injectivity-base-change-map-degree (blueprint/chap-Smooth-base-change/sec-Reduction-to-curves/lemma-injectivity-base-change-map-degree/_section.md)
\phantomsection\label{bp17}%
\begin{bpstatementit}{Lemma~8.3}{ (Injectivity of the base change map in degree $q$)}
Suppose \(q\geq1\) and \(\mathrm{BC}(g,n,q-1)\) holds. For every square obtained by base change as in the definition of \(\mathrm{BC}\) and every sheaf of \(\mathbf Z/n\)-modules \(M\) on \(Z_{\acute{e}t}\), the base change map \(b_M^q:g^{-1}R^qf_*M\to R^qf'_*g'^{-1}M\) in degree \(q\) is injective. An inclusion \(M\hookrightarrow N\) induces

\[
\operatorname{coker}(b_M^q)
\hookrightarrow\operatorname{coker}(b_N^q).
\]

For an embedding \(M\hookrightarrow I\) with \(I\) injective in the category of étale sheaves of \(\mathbf Z/n\)-modules, the cokernel of \(b_M^q\) embeds into \(R^qf'_*g'^{-1}I\).
\end{bpstatementit}

% proof: proof/lemma-injectivity-base-change-map-degree
\begin{proof}[{Proof}]\phantomsection\label{bp18}%
We embed \(M\) into an injective sheaf and compare the long exact sequences of the two sides. The induction hypothesis identifies everything in degree \(q-1\), and since the injective sheaf has no higher direct images upstairs, a diagram chase shows that \(b^q_M\) is injective with cokernel inside \(R^qf'_*g'^{-1}I\).

Choose \(0\to M\to I\to Q\to0\) with \(I\) as in the statement. By the long exact sequences, the rows of the following diagram are exact. Here the upper row ends in \(0\) because \(R^qf_*I=0\) for the injective sheaf \(I\), whereas the lower row does not, because \(g'^{-1}I\) need not be injective.

\[
\begin{CD}
g^{-1}R^{q-1}f_*I @>>> g^{-1}R^{q-1}f_*Q
@>>> g^{-1}R^qf_*M @>>>0\\
@V{\sim}VV @V{\sim}VV @VVV @VVV\\
R^{q-1}f'_*g'^{-1}I @>>> R^{q-1}f'_*g'^{-1}Q
@>>>R^qf'_*g'^{-1}M @>>>R^qf'_*g'^{-1}I.
\end{CD}
\]

The first two vertical arrows are isomorphisms by \(\mathrm{BC}(g,n,q-1)\), applied to \(I\) and to \(Q\). Exactness of the upper row says that \(g^{-1}R^qf_*M\) is the cokernel of the map \(g^{-1}R^{q-1}f_*I\to g^{-1}R^{q-1}f_*Q\), and exactness of the lower row says that the kernel of \(R^qf'_*g'^{-1}M\to R^qf'_*g'^{-1}I\) is the cokernel of the corresponding map in the lower row. Since the first two vertical isomorphisms identify these two maps, \(b_M^q\) is an isomorphism onto this kernel. In particular \(b^q_M\) is injective, and its cokernel embeds into \(R^qf'_*g'^{-1}I\), which is the last assertion. For \(M\subset N\), choose an injective \(I\) containing \(N\). What we have just shown, applied to \(M\subset I\) and to \(N\subset I\), embeds the cokernels of \(b_M^q\) and \(b_N^q\) compatibly into \(R^qf'_*g'^{-1}I\). As in [Stacks, \href{https://stacks.math.columbia.edu/tag/0F03}{0F03}], the map between the cokernels induced by \(M\subset N\) is compatible with these two embeddings and is therefore injective.
\end{proof}

% text: prose-08-08
\phantomsection\label{bp19}%
We use the lemma in the form of its second assertion: enlarging the coefficient sheaf does not lose a failure. In the next lemma this is how we pass from a scheme to an open subscheme or a closed complement.

% object: lemma-reduction-field (blueprint/chap-Smooth-base-change/sec-Reduction-to-curves/lemma-reduction-field/_section.md)
\phantomsection\label{bp20}%
\begin{bpstatementit}{Lemma~8.4}{ (Reduction to a field)}
Let \(q\geq1\). If \(\mathrm{BC}(g,n,q-1)\) holds but \(\mathrm{BC}(g,n,q)\) fails, then some failure occurs for an affine base change \(T\to Y\), direct image along a morphism \(\operatorname{Spec}\Omega\to T\) with \(\Omega\) a field, and a sheaf of \(\mathbf Z/n\)-modules on \(\operatorname{Spec}\Omega\).
\end{bpstatementit}

% proof: proof/lemma-reduction-field
\begin{proof}[{Proof}]\phantomsection\label{bp21}%
We shrink the source of the direct image, keeping a fixed nonzero element of the cokernel, until it becomes the spectrum of a field. A minimal closed subscheme carrying the failure exists by a limit argument, and its generic point still carries it. The one nonformal step is the third property below, which says that a failure on a scheme is seen either on an open subscheme or on its closed complement, and whose proof uses the injectivity of the base change map.

We first fix notation for the cokernels. Let \(T\) be a \(Y\)-scheme, put \(T'=Y'\times_YT\), and let \(g_T:T'\to T\) be the projection. For a qcqs morphism \(h:Z'\to T\) we write \(h':Z'_{T'}=Z'\times_TT'\to T'\) and \(g_{Z'}:Z'_{T'}\to Z'\) for its base changes, and for a sheaf \(N\) of \(\mathbf Z/n\)-modules on \(Z'_{\acute{e}t}\) we put

\[
b^q_{h,N}:g_T^{-1}R^qh_*N\longrightarrow R^qh'_*g_{Z'}^{-1}N,
\qquad
C(h,N)=\operatorname{coker}\bigl(b^q_{h,N}\bigr),
\]

a sheaf on \(T'_{\acute{e}t}\). These cokernels are functorial in two ways. A map of sheaves \(N\to N_1\) on \(Z'_{\acute{e}t}\) induces \(C(h,N)\to C(h,N_1)\), and a morphism \(u:Z''\to Z'\) induces \(C(h,N)\to C(hu,u^{-1}N)\), by pulling back along \(u\): upstairs through \(N\to u_*u^{-1}N\) and the edge map \(R^qh_*u_*\to R^q(hu)_*\), and downstairs in the same way. These pullback maps are compatible with the base change maps, and with composition in \(u\), because base change maps compose [Stacks, \href{https://stacks.math.columbia.edu/tag/0EZR}{0EZR}, \href{https://stacks.math.columbia.edu/tag/0EZS}{0EZS}].

Since \(\mathrm{BC}(g,n,q)\) fails, there are a \(Y\)-scheme \(T\), a qcqs morphism \(f:Z\to T\) and a sheaf \(M\) of \(\mathbf Z/n\)-modules on \(Z_{\acute{e}t}\) such that \(b^q_{f,M}\) is not an isomorphism. By Lemma \hyperref[bp17]{8.3} it is injective, so \(C(f,M)\neq0\), and we fix a geometric point \(\bar t'\) of \(T'\) and a nonzero element \(\xi\in C(f,M)_{\bar t'}\). For a qcqs morphism \(u:Z'\to Z\) let \(P(Z')\) mean that the image \(\xi_{Z'}\) of \(\xi\) in \(C(fu,u^{-1}M)_{\bar t'}\) is nonzero. The verification in [Stacks, \href{https://stacks.math.columbia.edu/tag/0F07}{0F07}] shows that \(P\) has three properties.

\begin{enumerate}
\item thickenings preserve \(P\), by topological invariance;
\item limits with affine transition maps preserve \(P\), by continuity;
\item if \(P(Z')\) holds and \(V\subset Z'\) is a quasi-compact open with closed complement \(W\), then \(P(V)\) or \(P(W)\) holds.
\end{enumerate}

To prove the third property, write \(h=fu:Z'\to T\) and \(M'=u^{-1}M\), and let \(j:V\to Z'\) and \(i:W\to Z'\) be the inclusions. We split \(M'\) into an open part and a closed part. Choose an injective sheaf \(J\) of \(\mathbf Z/n\)-modules on \(V_{\acute{e}t}\) containing \(j^{-1}M'\). Then the map

\[
M'\longrightarrow j_*J\oplus i_*i^{-1}M',
\]

whose components are adjoint to \(j^{-1}M'\subset J\) and to the identity of \(i^{-1}M'\), is injective. Indeed, at a geometric point of \(V\) its first component is the inclusion \(M'_{\bar v}\subset J_{\bar v}\), and at a geometric point of \(W\) its second component is the identity of \(M'_{\bar w}\). By the second assertion of Lemma \hyperref[bp17]{8.3}, the induced map

\[
C(h,M')\longrightarrow C(h,j_*J)\oplus C(h,i_*i^{-1}M')
\]

is injective as well. Since \(P(Z')\) holds, \(\xi_{Z'}\neq0\), and therefore its image in \(C(h,j_*J)_{\bar t'}\) or its image in \(C(h,i_*i^{-1}M')_{\bar t'}\) is nonzero. We show that the second case gives \(P(W)\) and the first gives \(P(V)\).

The closed summand is formal. The base change \(i':W_{T'}\to Z'_{T'}\) of \(i\) is again a closed immersion, and closed direct images are exact and commute with every base change (Lemma \bpout{4.27}). Hence, for a sheaf \(N\) on \(W_{\acute{e}t}\),

\[
R^qh_*(i_*N)=R^q(hi)_*N,
\qquad
R^qh'_*\bigl(g_{Z'}^{-1}i_*N\bigr)=R^qh'_*\bigl(i'_*g_W^{-1}N\bigr)
=R^q(hi)'_*g_W^{-1}N,
\]

and these identifications are compatible with the base change maps [Stacks, \href{https://stacks.math.columbia.edu/tag/0EZR}{0EZR}, \href{https://stacks.math.columbia.edu/tag/0EZS}{0EZS}]. Thus \(C(h,i_*N)=C(hi,N)\). For \(N=i^{-1}M'\) the image of \(\xi_{Z'}\) in \(C(h,i_*i^{-1}M')_{\bar t'}\) becomes \(\xi_W\), and so it is nonzero exactly when \(P(W)\) holds.

The open summand needs the induction hypothesis, because \(j\) is not finite. Let \(j':V_{T'}\to Z'_{T'}\) and \(g_V:V_{T'}\to V\) be the base changes of \(j\), and compare the Leray spectral sequences for \(hj=h\circ j\) and \(h'j'=h'\circ j'\). Upstairs, \(R^bj_*J=0\) for \(b>0\) because \(J\) is injective, so the Leray spectral sequence collapses and gives

\[
R^qh_*(j_*J)\xrightarrow{\ \sim\ }R^q(hj)_*J .
\]

Downstairs we use \(\mathrm{BC}(g,n,q-1)\). The scheme \(Z'\) is a \(Y\)-scheme through \(T\), and \(Z'_{T'}=Y'\times_YZ'\), so the square formed by \(j\), \(j'\), \(g_{Z'}\) and \(g_V\) is one of the squares in the definition of \(\mathrm{BC}\). Hence

\[
g_{Z'}^{-1}R^bj_*J\xrightarrow{\ \sim\ }R^bj'_*g_V^{-1}J
\qquad(0\le b\le q-1).
\]

For \(b=0\) this identifies \(g_{Z'}^{-1}j_*J\) with \(j'_*g_V^{-1}J\), and for \(0<b<q\) it shows that \(R^bj'_*g_V^{-1}J=0\). Consequently, in the Leray spectral sequence

\[
E_2^{a,b}=R^ah'_*R^bj'_*g_V^{-1}J\ \Longrightarrow\
R^{a+b}(h'j')_*g_V^{-1}J
\]

the rows \(0<b<q\) vanish. No differential can reach \(E^{q,0}\): the differential \(d_r\) into \(E_r^{q,0}\) starts at \(E_r^{q-r,r-1}\), which is zero for \(2\le r\le q\) because then \(0<r-1<q\), and for \(r>q\) because then \(q-r<0\). So \(E_2^{q,0}=E_\infty^{q,0}\), which is the bottom step of the filtration of the abutment, and the edge map

\[
e:R^qh'_*\bigl(j'_*g_V^{-1}J\bigr)\longrightarrow R^q(h'j')_*g_V^{-1}J
\]

is injective. Since base change maps compose [Stacks, \href{https://stacks.math.columbia.edu/tag/0EZR}{0EZR}, \href{https://stacks.math.columbia.edu/tag/0EZS}{0EZS}] and edge maps are natural, the square

\[
\begin{CD}
g_T^{-1}R^qh_*(j_*J) @>{b^q_{h,j_*J}}>> R^qh'_*\bigl(j'_*g_V^{-1}J\bigr)\\
@V{\wr}VV @VV{e}V\\
g_T^{-1}R^q(hj)_*J @>{b^q_{hj,J}}>> R^q(h'j')_*g_V^{-1}J
\end{CD}
\]

commutes. Its left arrow is an isomorphism and its right arrow is injective, so it induces an injective map \(C(h,j_*J)\to C(hj,J)\). (In the top row we have identified \(g_{Z'}^{-1}j_*J\) with \(j'_*g_V^{-1}J\) by the case \(b=0\); this does not change the cokernel.) The vertical arrows are the pullback maps along \(j\), so this map is pullback along \(j\) followed by \(j^{-1}j_*J=J\).

It remains to pass from \(J\) back to \(j^{-1}M'\). By naturality, the composite \(C(h,M')\to C(h,j_*J)\to C(hj,J)\) equals the composite \(C(h,M')\to C(hj,j^{-1}M')\to C(hj,J)\) of pullback along \(j\) and the map induced by \(j^{-1}M'\subset J\). In the first case the image of \(\xi_{Z'}\) in \(C(h,j_*J)_{\bar t'}\) is nonzero, hence so is its image in \(C(hj,J)_{\bar t'}\). This image is the image of \(\xi_V\) under the second composite, and so \(\xi_V\neq0\), which is \(P(V)\).

The field comes out of the three properties as follows. The intersection of a decreasing chain of closed subschemes is the limit of the chain, with affine transition maps. By the second property and Zorn's lemma there is therefore a minimal closed subscheme with \(P\). It is nonempty, since \(P\) fails for the empty scheme, and by the first property it is reduced. Let \(\xi\) be the generic point of an irreducible component. We claim that every affine open neighbourhood of \(\xi\) has \(P\). Indeed, by the third property, either the neighbourhood or its closed complement has \(P\), and the complement is a closed subscheme not containing \(\xi\), hence strictly smaller, so by minimality it does not have \(P\). Since the minimal closed subscheme is reduced, the limit of these neighbourhoods is \(\operatorname{Spec}\kappa(\xi)\). It has \(P\) by the second property (compare [Stacks, \href{https://stacks.math.columbia.edu/tag/0F06}{0F06}]), and so \(\kappa(\xi)\) is the field asserted in the lemma.
\end{proof}

% text: prose-08-09
\phantomsection\label{bp22}%
The field \(\Omega\) produced by this lemma is arbitrary, and the base \(T\) is still an arbitrary affine scheme. To bring them into the special situation of Lemma \hyperref[bp26]{8.6} we want to replace \(\Omega\) by its algebraic closure and \(T\) by a normal scheme. Both replacements are integral morphisms, and the next lemma shows that a failure survives them.

% object: lemma-integral-changes-source-base (blueprint/chap-Smooth-base-change/sec-Reduction-to-curves/lemma-integral-changes-source-base/_section.md)
\phantomsection\label{bp23}%
\begin{bpstatementit}{Lemma~8.5}{ (Integral changes of source and base)}
Assume \(\mathrm{BC}(g,n,q-1)\). A failure of base change in degree \(q\) persists after an integral surjective change of the source of the morphism along which we take direct image. It also persists after an integral change of the base through which this morphism factors.
\end{bpstatementit}

% proof: proof/lemma-integral-changes-source-base
\begin{proof}[{Proof}]\phantomsection\label{bp24}%
Both assertions are [Stacks, \href{https://stacks.math.columbia.edu/tag/0F04}{0F04}, \href{https://stacks.math.columbia.edu/tag/0F05}{0F05}]. They use three properties of direct image along an integral morphism \(\pi\): it is exact and conservative [Stacks, \href{https://stacks.math.columbia.edu/tag/04C2}{04C2}], even if \(\pi\) is not surjective, and it commutes with arbitrary base change [Stacks, \href{https://stacks.math.columbia.edu/tag/0EYP}{0EYP}]. By exactness, the Leray spectral sequence for a composite with \(\pi\) degenerates, which means that direct images along the composite are computed by applying \(\pi_*\) first.

We use the notation \(b^q_{h,N}\) and \(C(h,N)=\operatorname{coker}b^q_{h,N}\) of the proof of Lemma \hyperref[bp20]{8.4}, for a \(Y\)-scheme \(T\), its base change \(g_T:T'=Y'\times_YT\to T\), a qcqs morphism \(h:Z\to T\) and a sheaf \(N\) of \(\mathbf Z/n\)-modules on \(Z_{\acute{e}t}\). By Lemma \hyperref[bp17]{8.3}, \(b^q_{h,N}\) is injective, so a failure of base change in degree \(q\) means \(C(h,N)\neq0\).

For the first assertion, let \(f:Z\to T\) and \(M\) give a failure, so \(C(f,M)\neq0\), and let \(\pi:\widetilde Z\to Z\) be integral and surjective. We show that \(C(f\pi,\pi^{-1}M)\neq0\). Since \(\pi\) is surjective, the unit \(M\to\pi_*\pi^{-1}M\) is injective, as one sees on stalks. By the second assertion of Lemma \hyperref[bp17]{8.3} it induces an injection \(C(f,M)\hookrightarrow C(f,\pi_*\pi^{-1}M)\), and so \(C(f,\pi_*\pi^{-1}M)\neq0\). It remains to identify \(C(f,\pi_*\pi^{-1}M)\) with \(C(f\pi,\pi^{-1}M)\). Let \(\pi':\widetilde Z_{T'}\to Z_{T'}\) be the base change of \(\pi\), again integral. Upstairs, \(\pi_*\) is exact, so the Leray spectral sequence for \(f\circ\pi\) degenerates and gives \(R^qf_*(\pi_*\pi^{-1}M)=R^q(f\pi)_*\pi^{-1}M\). Downstairs, \(\pi_*\) commutes with base change, \(g_Z^{-1}\pi_*\pi^{-1}M\simeq\pi'_*g_{\widetilde Z}^{-1}\pi^{-1}M\), and \(\pi'_*\) is exact, so in the same way

\[
R^qf'_*\bigl(g_Z^{-1}\pi_*\pi^{-1}M\bigr)
\simeq R^qf'_*\pi'_*\bigl(g_{\widetilde Z}^{-1}\pi^{-1}M\bigr)
=R^q(f\pi)'_*\bigl(g_{\widetilde Z}^{-1}\pi^{-1}M\bigr).
\]

These identifications are compatible with the base change maps [Stacks, \href{https://stacks.math.columbia.edu/tag/0EZR}{0EZR}, \href{https://stacks.math.columbia.edu/tag/0EZS}{0EZS}], so they carry \(b^q_{f,\pi_*\pi^{-1}M}\) to \(b^q_{f\pi,\pi^{-1}M}\). Hence \(C(f\pi,\pi^{-1}M)\simeq C(f,\pi_*\pi^{-1}M)\neq0\), and base change fails for the new source.

For the second assertion, let \(f:Z\to T\) factor as \(f=\pi\circ f_1\) with \(\pi:T_1\to T\) integral and \(f_1:Z\to T_1\), and suppose base change fails for \(f\) and \(M\), that is, \(C(f,M)\neq0\). We show that it fails for \(f_1\) and \(M\) over the new base \(T_1\). Put \(T'_1=Y'\times_YT_1=T'\times_TT_1\), and let \(\pi':T'_1\to T'\) be the base change of \(\pi\), again integral. The base change map for \(f_1\) is

\[
b^q_{f_1,M}:g_{T_1}^{-1}R^qf_{1*}M\longrightarrow R^qf'_{1*}g_Z^{-1}M
\qquad\text{on }T'_{1,\acute{e}t}.
\]

We push it down to \(T'\) by \(\pi'_*\). Since \(\pi'_*\) is exact and \(\pi_*\) commutes with base change, \(\pi'_*g_{T_1}^{-1}R^qf_{1*}M\simeq g_T^{-1}\pi_*R^qf_{1*}M
=g_T^{-1}R^qf_*M\), where the last equality is the degenerate Leray spectral sequence for \(f=\pi f_1\). In the same way \(\pi'_*R^qf'_{1*}g_Z^{-1}M=R^qf'_*g_Z^{-1}M\). Under these identifications \(\pi'_*\bigl(b^q_{f_1,M}\bigr)=b^q_{f,M}\), again because base change maps compose. Since \(\pi'_*\) is exact, it follows that \(\pi'_*C(f_1,M)\simeq C(f,M)\neq0\), so \(C(f_1,M)\neq0\), which is the assertion.
\end{proof}

% text: prose-08-10
\phantomsection\label{bp25}%
With the two previous lemmas we may take the source of the direct image to be the spectrum of an algebraically closed field. Its cohomology vanishes in positive degrees, so the source of the base change map is zero and the failure sits in the target, a direct image along an extension of the ground field.

% object: lemma-reduction-curve-over-normal-base (blueprint/chap-Smooth-base-change/sec-Reduction-to-curves/lemma-reduction-curve-over-normal-base/_section.md)
\phantomsection\label{bp26}%
\begin{bpstatementit}{Lemma~8.6}{ (Reduction to a curve over a normal base)}
Suppose \(g:Y'\to Y\) is affine and smooth of relative dimension one, \(n\) is invertible on \(Y\), and \(\mathrm{BC}(g,n,q-1)\) holds. If \(\mathrm{BC}(g,n,q)\) fails, there are an affine scheme \(S=\operatorname{Spec}R\) over \(Y\) with \(R\) a normal domain whose fraction field \(L\) is algebraically closed, an algebraically closed extension field \(\Omega\) of \(L\), an integer \(m\mid n\), and, for the smooth affine curve \(C=Y'\times_YS\), the cartesian square

\[
\begin{CD}
C_\Omega @>{h}>> C\\
@VVV @VVV\\
\operatorname{Spec}\Omega @>>> S
\end{CD}
\]

such that

\[
R^qh_*\underline{\mathbf Z/m}\ne0.
\]
\end{bpstatementit}

% text: prose-08-11
\phantomsection\label{bp27}%
Here \(\underline{\mathbf Z/m}\) is the classical constant sheaf.

% proof: proof/lemma-reduction-curve-over-normal-base
\begin{proof}[{Proof of Lemma 8.6}]\phantomsection\label{bp28}%
We first simplify the source of the direct image and the coefficients, then show that the source of the base change map vanishes, and finally replace the base by a normal one.

By Lemma \hyperref[bp20]{8.4} there is a failure with direct image from the spectrum of a field \(\Omega\). Since the map from the spectrum of an algebraic closure of \(\Omega\) to \(\operatorname{Spec}\Omega\) is integral and surjective, Lemma \hyperref[bp23]{8.5} allows us to replace \(\Omega\) by an algebraic closure. On this geometric point, a sheaf of \(\mathbf Z/n\)-modules is a \(\mathbf Z/n\)-module, hence a direct sum of cyclic groups \(\mathbf Z/m\) with \(m\mid n\). As the base change map commutes with this direct sum by continuity of qcqs direct images, base change still fails for some summand \(\mathbf Z/m\).

The source of the base change map is the pullback of \(R^qf_*\underline{\mathbf Z/m}\) for \(f:\operatorname{Spec}\Omega\to T\), and we show that it is zero. The stalk of this sheaf at a geometric point \(\bar t\) of \(T\) is the cohomology of \(\operatorname{Spec}\Omega\times_T\operatorname{Spec}\mathcal O^{\mathrm{sh}}_{T,\bar t}\). This scheme is a limit of étale \(\Omega\)-schemes, which, since \(\Omega\) is algebraically closed, are disjoint unions of copies of \(\operatorname{Spec}\Omega\). By continuity of cohomology the stalk therefore vanishes for \(q\geq1\). As the source of the base change map is zero, the failure means that the target \(R^qh_*\underline{\mathbf Z/m}\) is nonzero.

The extension theorem for covers will require a normal base with algebraically closed fraction field. If the base is \(\operatorname{Spec}B\), we replace it by \(\operatorname{Spec}R\), where \(R\) is the integral closure in \(\Omega\) of the image of \(B\). This is allowed, because the morphism \(\operatorname{Spec}\Omega\to\operatorname{Spec}B\) factors through \(\operatorname{Spec}R\), and \(\operatorname{Spec}R\to\operatorname{Spec}B\) is integral, so the failure persists by Lemma \hyperref[bp23]{8.5}. Moreover, the fraction field \(L\) of \(R\) is the algebraic closure of \(\operatorname{Frac}(\operatorname{im}(B\to \Omega))\) in \(\Omega\), hence algebraically closed, although in general \(\Omega/L\) is transcendental. We are now in the situation of [Stacks, \href{https://stacks.math.columbia.edu/tag/0F08}{0F08}].
\end{proof}

% object: example-reduction-end-affine-line (blueprint/chap-Smooth-base-change/sec-Reduction-to-curves/example-reduction-end-affine-line/_section.md)
\phantomsection\label{bp29}%
\begin{bpstatement}{Example~8.7}{ (Where the reduction ends)}
The square of the lemma is itself an instance of smooth base change, and the simplest case shows what remains to be proved. Let \(Y=S=\operatorname{Spec}L\) with \(L\) algebraically closed, let \(g:Y'=C=\mathbf A^1_L\to Y\) be the structure map, and take direct image along \(f:X=\operatorname{Spec}\Omega\to Y\) for an algebraically closed extension \(\Omega\) of \(L\), so that \(X'=\mathbf A^1_\Omega\) and \(f'=h\). The sheaf \(R^qf_*\mathbf Z/m\) on \(\operatorname{Spec}L\) is \(H^q(\operatorname{Spec}\Omega,\mathbf Z/m)\), which vanishes for \(q\geq1\). So smooth base change for this square says exactly that \(R^qh_*\mathbf Z/m=0\) for \(q\geq1\). For \(m\) invertible in \(L\) this is the content of Proposition \hyperref[bp32]{8.8} below.

Let \(q=1\) and \(m\) invertible in \(L\). The sections over \(C\) itself are no problem: choosing \(\mu_m\simeq\mathbf Z/m\), the Kummer sequence gives \(H^1(\mathbf A^1_\Omega,\mathbf Z/m)=\Omega^*/(\Omega^*)^m=0\), since \(\Gamma(\mathbf A^1_\Omega,\mathcal O)^*=\Omega^*\) is divisible and \(\operatorname{Pic}(\mathbf A^1_\Omega)=0\). However, the presheaf whose sheafification is \(R^1h_*\mathbf Z/m\) has other sections. For example, on the open \(U=\mathbf G_{m,L}\) of \(C\) the group \(H^1(\mathbf G_{m,\Omega},\mathbf Z/m)\) is cyclic of order \(m\), and such a class has to be killed by an étale covering of \(U\), as we shall see after the statement of the proposition.

For \(m=p\) in characteristic \(p\) the argument breaks down already over \(C\) itself. Indeed, the group \(H^1(\mathbf A^1_\Omega,\mathbf F_p)\) is the cokernel of \(u\mapsto u^p-u\) on \(\Omega[x]\), and for \(\omega\in\Omega\) transcendental over \(L\) it contains the class of \(\omega x\), which does not come from \(C\) (compare the class of \(tx\) at the beginning of the chapter).
\end{bpstatement}


% object: sec-Change-of-base-field (blueprint/chap-Smooth-base-change/sec-Change-of-base-field/_section.md), content hidden
\subsection*{\texorpdfstring{8.6\quad Change of base field}{8.6 Change of base field}}
\phantomsection\label{bp30}%
\addcontentsline{toc}{subsection}{\texorpdfstring{8.6\quad Change of base field}{8.6 Change of base field}}
% text: prose-08-12
\phantomsection\label{bp31}%
Lemma \hyperref[bp26]{8.6} turned a failure of \(\mathrm{BC}(g,n,q)\) into a single nonvanishing sheaf \(R^qh_*\mathbf Z/m\ne0\), where \(h:C_\Omega\to C\) is the base change of a smooth affine curve \(C\) over a normal base \(S\) to a larger algebraically closed field \(\Omega\). Proposition \hyperref[bp32]{8.8} below shows that this sheaf vanishes, which contradicts the lemma and finishes the induction. As the sheaf \(R^qh_*\mathbf Z/m\) is the sheafification of \(U\mapsto H^q(U_\Omega,\mathbf Z/m)\) on \(C_{\acute{e}t}\), we must show that classes on \(U_\Omega\) die after étale localization on \(U\). In degrees \(q\geq2\) there is nothing to kill, since affine curves over an algebraically closed field have cohomological dimension one.

In degree one the problem is about covers. Recall that \(S=\operatorname{Spec}R\) with \(R\) a normal domain whose fraction field \(L\) is algebraically closed, and \(L\subset\Omega\) (Lemma \hyperref[bp26]{8.6}). Over the generic point of \(S\), a class on \(C_\Omega\) already comes from \(C_L\), by the invariance of \(H^1\) of a curve under extension of algebraically closed fields (Theorem \bpout{4.20}). In other words, it is given by a \(\mathbf Z/m\)-torsor over \(C_L\), so it suffices to find an étale cover of \(C\) on which this torsor becomes trivial. To find one, we take a connected component of the torsor and extend it over \(C\) by the extension theorem for Galois covers of order invertible on \(S\). This theorem needs \(S\) normal and a section of \(C\to S\). The section is available: since \(L\) is algebraically closed, the generic fibre \(C_L\) has a rational point, and after étale localization on \(S\) this point extends to a section.

We use that \(m\) is invertible in two places, and both fail for \(p\)-torsion. First, the invariance of \(H^1\) fails: the Artin–Schreier class of \(tx\) on \(\mathbf A^1_{k(t)^{\mathrm{sep}}}\) does not come from \(\mathbf A^1_k\). Second, the extension theorem can fail for covers of degree divisible by a residue characteristic, because such covers may be wildly ramified along the special fibres.

% object: proposition-base-change-larger-algebraically-closed (blueprint/chap-Smooth-base-change/sec-Change-of-base-field/proposition-base-change-larger-algebraically-closed/_section.md)
\phantomsection\label{bp32}%
\begin{bpstatementit}{Proposition~8.8}{ (Base change to a larger algebraically closed field)}
Let \(S=\operatorname{Spec}R\) for a normal domain \(R\) with algebraically closed fraction field \(L\), and let \(a:C\to S\) be affine and smooth of relative dimension one. Let \(\Omega/L\) be an extension of algebraically closed fields, let \(m\) be a positive integer invertible on \(S\), and let \(h:C_\Omega\to C\) be the projection. Then

\[
R^qh_*\underline{\mathbf Z/m}=0\qquad(q\geq1)
\]

on \(C_{\mathrm{pro}\acute{e}t}\).
\end{bpstatementit}

% text: prose-08-13
\phantomsection\label{bp33}%
The simplest case already shows how the proof goes. Take \(C=\mathbf G_{m,S}\) with coordinate \(x\). Then \(H^1(\mathbf G_{m,\Omega},\mathbf Z/m)\) is cyclic of order \(m\), generated by the class of the Kummer cover \(y^m=x\). Here we view this cover as a \(\mathbf Z/m\)-torsor through a primitive \(m\)-th root of unity, which lies in \(R\) because \(R\) is normal and \(L\) is algebraically closed. The point is that the cover is already defined over \(\mathbf G_{m,S}\), and it is finite étale there because \(m\) is invertible on \(S\), so the class becomes trivial after pullback to it. The proof below produces such a cover for an arbitrary class.

% proof: proof/proposition-base-change-larger-algebraically-closed
\begin{proof}[{Proof of Proposition 8.8}]\phantomsection\label{bp34}%
We describe \(R^qh_*\) as a sheafification and show that each local class dies after a suitable étale covering. In degrees at least two this is the cohomological dimension of affine curves. In degree one we first descend the class to \(C_L\), then trivialize it on a finite étale Galois cover of \(C_L\), and finally extend this cover over \(C\).

\emph{Step 1: reduction to local classes.} We may work on the étale site, because \(h\) is qcqs and direct images can be compared (Theorem \bpout{2.38}). There, [Stacks, \href{https://stacks.math.columbia.edu/tag/03Q8}{03Q8}] identifies \(R^qh_*\underline{\mathbf Z/m}\) with the sheafification of the presheaf

\[
U\longmapsto H^q(U\times_S\operatorname{Spec}\Omega,
\mathbf Z/m),
\qquad U\to C\text{ \'etale},
\]

so we must show that every such class vanishes on an étale covering of \(U\).

\emph{Step 2: degrees \(q\geq2\).} In these degrees the groups themselves vanish for affine \(U\), which we may assume because the question is local on \(U\). Indeed, \(U_\Omega\) is then affine and étale over \(C_\Omega\). Hence it is empty or a finite disjoint union of smooth affine curves over the algebraically closed field \(\Omega\), and \((C_2)\) of Theorem \bpout{4.20} kills its cohomology in degrees \(q\geq2\).

\emph{Step 3: a section and connected fibres.} From now on \(q=1\). The extension theorem we want to use requires a section and geometrically connected fibres. The map \(U\to S\) need not have them, but we may arrange them locally. By [Stacks, \href{https://stacks.math.columbia.edu/tag/0EY4}{0EY4}], after an étale covering of \(U\) the map \(U\to S\) factors as \(U\to V\to S\) with \(U\) and \(V\) affine, \(V\to S\) étale, and \(U\to V\) smooth with geometrically connected fibres and a section. Moreover, the separated étale map \(V\to S\) is a finite disjoint union of open immersions [Stacks, \href{https://stacks.math.columbia.edu/tag/09Z9}{09Z9}], because \(S\) is normal and integral with separably closed function field. As we only have to kill classes on \(U_\Omega\), we may therefore replace \(S\) by one of these opens and \(C\) by the part of \(U\) over it, which does not change the fraction field \(L\). Afterwards there is a section \(s:S\to C\) and the fibres are geometrically connected.

\emph{Step 4: descent to \(C_L\) and a trivializing cover.} Let \(\xi\in H^1(C_\Omega,\mathbf Z/m)\). It is the pullback of a unique class \(\xi_L\in H^1(C_L,\mathbf Z/m)\), because the pullback \(H^1(C_L,\mathbf Z/m)\to H^1(C_\Omega,\mathbf Z/m)\) is an isomorphism by \((C_6)\) of Theorem \bpout{4.20}, applied to the affine curve \(C_L\) as in [Stacks, \href{https://stacks.math.columbia.edu/tag/03SC}{03SC}]; this step uses that \(m\) is prime to the characteristic. Next, the classification of \(H^1\) by torsors [Stacks, \href{https://stacks.math.columbia.edu/tag/03AJ}{03AJ}] and Lemma \bpout{2.41} show that \(\xi_L\) is the class of a finite étale \(\mathbf Z/m\)-torsor \(P\to C_L\). Choose a connected component \(Z\) of \(P\) and let \(G\subseteq\mathbf Z/m\) be its stabilizer. As the curve \(C_L\) is connected, \(Z\to C_L\) is a finite étale Galois cover with group \(G\). The class \(\xi_L\) vanishes on \(Z\), because the pullback of \(P\) to \(Z\) has the diagonal section.

\emph{Step 5: extending the cover over \(C\).} We use the extension theorem for generic Galois covers [Stacks, \href{https://stacks.math.columbia.edu/tag/0EZJ}{0EZJ}]: let \(S\) be integral, normal and qcqs with generic point \(\eta\), and let \(C\to S\) be qcqs and smooth with geometrically connected fibres and a section \(s\). If \(Z\to C_\eta\) is a finite étale Galois cover with group \(G\) of order invertible on \(S\), and \(Z\) has a \(\kappa(\eta)\)-rational point over \(s(\eta)\), then \(Z\) extends to a finite étale Galois cover \(\widetilde Z\to C\) with group \(G\). The idea is that a cover of order prime to the residue characteristics ramifies only tamely, and splitting over the section then rules out ramification along the irreducible special fibres, and at points of codimension at least two the cover extends by purity of the branch locus.

The hypotheses hold with \(C_\eta=C_L\). Indeed, the order \(|G|\) divides \(m\) and is therefore invertible on \(S\). Moreover, the fibre of \(Z\) over \(s(\eta)\) is nonempty, as \(Z\to C_L\) is finite étale and surjective, and it has an \(L\)-rational point because \(L\) is algebraically closed. So \(\widetilde Z\) exists.

We check that \(\widetilde Z\to C\) is an étale covering, that is, that it is surjective. The scheme \(C\) is connected, since \(C\to S\) is open and surjective with connected fibres over a connected base. Hence the image of \(\widetilde Z\) in \(C\), which is open and closed and contains the generic fibre, is all of \(C\). Finally, \(\xi\) dies on this covering. Indeed, \(\widetilde Z\times_C C_\Omega=Z\times_{C_L}C_\Omega\), so the pullback of \(\xi\) to \(\widetilde Z\times_CC_\Omega\) is the pullback of \(\xi_L|_Z=0\).
\end{proof}

% proof: proof/theorem-smooth-base-change-classical-torsion
\begin{proof}[{Proof of Theorem 8.1}]\phantomsection\label{bp35}%
We argue by induction on \(q\). The degree-zero calculation starts the induction. By the reductions above, a hypothetical failure in degree \(q\) gives a nonzero sheaf \(R^qh_*\underline{\mathbf Z/m}\), and Proposition \hyperref[bp32]{8.8} says that this sheaf is zero.

A smooth morphism is flat, locally of finite presentation, and has geometrically reduced fibres, and so the case of degree zero is Lemma \hyperref[bp11]{8.2}. Let \(q\geq1\) and assume the theorem in all degrees less than \(q\). Fix \(n\) invertible on \(Y\). Since smoothness is stable under base change, we may apply the induction hypothesis after every base change \(T\to Y\), which shows that \(\mathrm{BC}(g,n,q-1)\) holds. By étale coordinates and composition of base change maps, it suffices to treat the case where \(g\) is affine and smooth of relative dimension one.

If \(\mathrm{BC}(g,n,q)\) failed, Lemma \hyperref[bp26]{8.6} would give \(S\), \(C\), \(h\) and \(m\) with \(R^qh_*\underline{\mathbf Z/m}\neq0\). These satisfy the hypotheses of Proposition \hyperref[bp32]{8.8}: by construction \(S\) is the spectrum of a normal domain with algebraically closed fraction field, \(m\) is invertible on \(S\) because it divides \(n\), and \(C\) is a smooth affine curve over \(S\). Therefore the conclusion \(R^qh_*\underline{\mathbf Z/m}=0\) of the proposition is a contradiction. Hence \(\mathrm{BC}(g,n,q)\) holds for affine \(g\) of relative dimension one, and by the reduction above it holds for every smooth \(g\).

Since \(f\) and \(f'\) are qcqs, both sides commute with the filtered union \(F=\bigcup_nF[n]\), which gives the theorem for \(F\) on the étale sites. Finally, Theorem \bpout{2.38} transfers the isomorphism and the base change map from the étale sites to \(Y'_{\mathrm{pro}\acute{e}t}\).
\end{proof}


% object: sec-Bounded-below-complexes (blueprint/chap-Smooth-base-change/sec-Bounded-below-complexes/_section.md), content hidden
\subsection*{\texorpdfstring{8.7\quad Bounded below complexes}{8.7 Bounded below complexes}}
\phantomsection\label{bp36}%
\addcontentsline{toc}{subsection}{\texorpdfstring{8.7\quad Bounded below complexes}{8.7 Bounded below complexes}}
% text: prose-08-14
\phantomsection\label{bp37}%
Theorem \hyperref[bp7]{8.1} concerns a single torsion sheaf, but its applications concern complexes. Smooth base change on \(D^b_{\mathrm{ctf}}\) (Theorem \hyperref[bp45]{8.11}) is reduced, for integral coefficients, to complexes with finite coefficients, and the Künneth formula (Section \hyperref[bp49]{8.9}) is a statement about complexes. So we need smooth base change for complexes of torsion sheaves, and bounded below complexes suffice, since objects of \(D^b_{\mathrm{ctf}}\) are bounded.

% object: theorem-smooth-base-change-classical-torsion-2 (blueprint/chap-Smooth-base-change/sec-Bounded-below-complexes/theorem-smooth-base-change-classical-torsion-2/_section.md)
\phantomsection\label{bp38}%
\begin{bpstatementit}{Theorem~8.9}{ (Smooth base change for classical torsion complexes)}
In (8.1), let \(g\) be smooth and \(f\) quasi-compact and quasi-separated. Let \(\Lambda\) be a discrete commutative ring annihilated by an integer \(n\) invertible on \(Y\). For every

\[
K\in D^+_{\mathrm{class}}(X_{\mathrm{pro}\acute{e}t},\Lambda)
\]

the canonical map (4.5) is an isomorphism.
\end{bpstatementit}

% proof: proof/theorem-smooth-base-change-classical-torsion-2
\begin{proof}[{Proof}]\phantomsection\label{bp39}%
We first treat a single sheaf, where the only point is \(\Lambda\)-linearity, and then pass to complexes by comparing hypercohomology spectral sequences.

For a classical sheaf the assertion is Theorem \hyperref[bp7]{8.1}, applied to the underlying abelian sheaf, which is killed by \(n\). Since the base change map is constructed from units and counits, it is \(\Lambda\)-linear, and so it is an isomorphism of \(\Lambda\)-module sheaves.

For a complex, we first note that by Theorem \bpout{2.38} and Corollary \bpout{2.37} the direct and inverse images involved are classical, and that Lemma \bpout{4.7} identifies the comparison with (4.5). So we may argue on the étale sites. Choose an integer \(m\) with \(K\in D^{\geq m}\). Since both \(g^*\) and \(g'^*\) are exact, they commute with taking cohomology sheaves, and the two sides of (4.5) are computed by the hypercohomology spectral sequences for \(Rf_*\) and \(Rf'_*\), applied to \(K\) and to \(g'^*K\). These spectral sequences have \(E_2\) terms

\[
g^*R^pf_*H^q(K),\qquad R^pf'_*g'^*H^q(K),
\qquad p\geq0,\quad q\geq m,
\]

and the base change map induces on them the comparison for sheaves, which is an isomorphism by the preceding paragraph. Since in each total degree only finitely many pairs \((p,q)\) occur, these shifted first quadrant spectral sequences converge; this is where we use that \(K\) is bounded below. A map of convergent spectral sequences that is an isomorphism on \(E_2\) is an isomorphism on abutments, so the base change map induces isomorphisms \(g^*R^if_*K\to R^if'_*g'^*K\) for all \(i\), and so it is an isomorphism in the derived category.
\end{proof}

% text: prose-08-15
\phantomsection\label{bp40}%
For the theorem on \(D^b_{\mathrm{ctf}}\) we only need the case of a finite admissible ring \(A\), where the hypothesis holds with \(n=\ell^r\) for \(\ell^rA=0\).


% object: sec-ell-adic-coefficients-3 (blueprint/chap-Smooth-base-change/sec-ell-adic-coefficients-3/_section.md), content hidden
\subsection*{\texorpdfstring{8.8\quad $\ell$-adic coefficients}{8.8 ell-adic coefficients}}
\phantomsection\label{bp41}%
\addcontentsline{toc}{subsection}{\texorpdfstring{8.8\quad $\ell$-adic coefficients}{8.8 ell-adic coefficients}}
% text: prose-08-16
\phantomsection\label{bp42}%
For every admissible coefficient ring, smooth base change takes the following form. For finite rings it is a case of smooth base change for torsion complexes (Theorem \hyperref[bp38]{8.9}), and we deduce the other rings by Proposition \bpout{7.12}.

On the pro-étale site the base change map for \(\mathcal O_E\)-coefficients is a single map of complexes, defined by the same adjunction as for finite rings. Lemma \hyperref[bp43]{8.10} identifies its reduction modulo \(\pi_E\) with the base change map for \(\mathcal O_E/\pi_E\), so derived Nakayama gives the integral case from the torsion theorem. The rational cases then follow by inverting \(\ell\) and by descent to a finite extension.

% object: lemma-change-coefficients (blueprint/chap-Smooth-base-change/sec-ell-adic-coefficients-3/lemma-change-coefficients/_section.md)
\phantomsection\label{bp43}%
\begin{bpstatementit}{Lemma~8.10}{ (Change of coefficients)}
Suppose that the schemes in (8.1) are of finite type over \(k\). Let \(\Lambda\to\Lambda'\) be a coefficient map of Proposition \bpout{7.10} and let \(K\in D^b_{\mathrm{ctf}}(X,\Lambda)\). Then the square

\[
\begin{CD}
(g^*Rf_*K)_{\Lambda'} @>>> (Rf'_*g'^*K)_{\Lambda'}\\
@VVV @VVV\\
g^*Rf_*(K_{\Lambda'}) @>>> Rf'_*g'^*(K_{\Lambda'})
\end{CD}
\]

commutes. Here the upper arrow is the extension of scalars of the base change map (4.5) for \(K\), the lower arrow is the base change map for \(K_{\Lambda'}\), and the vertical arrows are the composites of the change-of-coefficients isomorphisms of Proposition \bpout{7.10} for \(g^*\) and \(Rf_*\), respectively for \(Rf'_*\) and \(g'^*\).
\end{bpstatementit}

% proof: proof/lemma-change-coefficients
\begin{proof}[{Proof}]\phantomsection\label{bp44}%
The base change map is built from units, counits and composition isomorphisms for inverse images. By Proposition \bpout{7.10} the change-of-coefficients maps are compatible with all of them, and so the square commutes.
\end{proof}

% object: theorem-smooth-base-change-all-admissible (blueprint/chap-Smooth-base-change/sec-ell-adic-coefficients-3/theorem-smooth-base-change-all-admissible/_section.md)
\phantomsection\label{bp45}%
\begin{bpstatementit}{Theorem~8.11}{ (Smooth base change for all admissible coefficients)}
In (8.1), assume that all schemes are of finite type over \(k\) and that \(g\) is smooth. For every admissible coefficient ring

\[
\Lambda\in\{A,\mathbf Z_\ell,\mathcal O_E,\mathbf Q_\ell,E,\overline{\mathbf Q}_\ell\}
\]

and every \(K\in D^b_{\mathrm{ctf}}(X,\Lambda)\), the base change map (4.5) is an isomorphism

\[
g^*Rf_*K\xrightarrow{\sim}Rf'_*g'^*K.
\tag{8.2}
\]
\end{bpstatementit}

% text: prose-08-17
\phantomsection\label{bp46}%
For integral and rational coefficients we have \(D^b_{\mathrm{ctf}}=\Dbc\), so in those cases the theorem concerns all bounded constructible complexes. The isomorphisms (8.2) respect composition of cartesian squares and change of coefficients, because the base change map itself does (Lemmas \bpout{4.7} and \hyperref[bp43]{Lemma 8.10}).

% proof: proof/theorem-smooth-base-change-all-admissible
\begin{proof}[{Proof of Theorem 8.11}]\phantomsection\label{bp47}%
For finite rings the theorem is the torsion theorem once we know that \(K\) is classical, and the other rings follow by the reduction principle. By Proposition \bpout{7.1}, both sides lie in \(D^b_{\mathrm{ctf}}(Y',\Lambda)\), so we only have to show that the canonical map (4.5) is invertible.

\emph{Finite rings \(A\).} The torsion theorem is stated for classical complexes, so we first check that \(K\) is one. By Theorem \bpout{2.43}, constructible pro-étale sheaves with discrete Noetherian coefficients are classical, and since the cohomology sheaves of \(K\) are constructible, \(K\) is a bounded classical complex. Since \(A\) is killed by a power of \(\ell\), which is invertible on \(Y\), smooth base change for torsion complexes (Theorem \hyperref[bp38]{8.9}) applies to \(K\).

\emph{Other rings.} We apply Proposition \bpout{7.12} with \(F_\Lambda=g^*Rf_*\), \(G_\Lambda=Rf'_*g'^*\) and \(\theta_\Lambda\) the base change map. Both functors are exact, as the proposition requires, and the base change map commutes with change of coefficients, which is Lemma \hyperref[bp43]{8.10}. Since \(\theta_A\) is an isomorphism for every finite \(A\) by the first case, the proposition shows that \(\theta_\Lambda\) is an isomorphism for every admissible \(\Lambda\).
\end{proof}

% object: remark-smooth-base-change-proper-base (blueprint/chap-Smooth-base-change/sec-ell-adic-coefficients-3/remark-smooth-base-change-proper-base/_section.md)
\phantomsection\label{bp48}%
\begin{bpstatement}{Remark~8.12}{ (Smooth base change, proper base change and base change for $f^!$)}
In (8.1), smooth base change assumes that \(g\) is smooth and allows any \(f\) of finite type, whereas proper base change assumes that \(f\) is proper and allows any \(g\). Smooth purity (Chapter \bpout{9}) gives a third base change isomorphism \(g'^*f^!\simeq f'^!g^*\), for the exceptional inverse image, when \(f\) is smooth and \(g\) is arbitrary. By contrast, neither smooth nor proper base change applies to the open immersion \(j:\mathbf G_m\hookrightarrow\mathbf A^1\) of the beginning of this chapter, where base change along the immersion of the origin fails: Theorem \hyperref[bp45]{8.11} does not apply because the immersion of the origin is not smooth, and proper base change does not apply because \(j\) is not proper.
\end{bpstatement}


% object: sec-The-Kunneth-formula (blueprint/chap-Smooth-base-change/sec-The-Kunneth-formula/_section.md), content hidden
\subsection*{\texorpdfstring{8.9\quad The Künneth formula}{8.9 The Künneth formula}}
\phantomsection\label{bp49}%
\addcontentsline{toc}{subsection}{\texorpdfstring{8.9\quad The Künneth formula}{8.9 The Künneth formula}}
% text: prose-08-18
\phantomsection\label{bp50}%
The trace formula needs the Künneth formula for ordinary direct images (Proposition \hyperref[bp64]{8.18}), which is a base change statement for \(Rf_*\). Base change for \(Rf_*\) fails in general, as the open immersion \(j:\mathbf G_m\hookrightarrow\mathbf A^1\) of the beginning of the chapter shows. Along a product projection \(Y\times T\to Y\), however, base change holds (Lemma \hyperref[bp51]{8.13}). To prove this we compactify, after which only an open immersion remains, and for an open immersion the statement comes down to two results: base change along extensions of the ground field, a consequence of smooth base change, and the Künneth formula with one proper factor, a consequence of proper base change. In this section \(A\) is a finite admissible ring, and we write \(M\boxtimes N=p_1^*M\otimes^{\mathbf L}p_2^*N\) for the external tensor product.

We start with the base change statement, which is the case of the Künneth formula in which the second coefficient is constant.

% object: lemma-base-change-along-product-projection (blueprint/chap-Smooth-base-change/sec-The-Kunneth-formula/lemma-base-change-along-product-projection/_section.md)
\phantomsection\label{bp51}%
\begin{bpstatementit}{Lemma~8.13}{ (Base change along a product projection)}
Let \(f:X\to Y\) be a morphism of separated \(k\)-schemes of finite type, let \(T\) be a separated \(k\)-scheme of finite type, and let \(M\in\Dbc(X,A)\).  Let \(p_X:X\times T\to X\) and \(p_Y:Y\times T\to Y\) be the projections.  Then the base change map

\[
p_Y^*Rf_*M\longrightarrow R(f\times\operatorname{id}_T)_*p_X^*M
\]

is an isomorphism.
\end{bpstatementit}

% proof: proof/lemma-base-change-along-product-projection
\begin{proof}[{Proof}]\phantomsection\label{bp52}%
The idea is to compactify \(f\), so that only an open immersion remains, and for an open immersion base change along a product projection holds even when \(T\) is singular. We argue on the étale site, where \(M\) is a bounded complex with torsion coefficients, and transport the result to the pro-étale site by Theorem \bpout{2.38}.

Choose a compactification \(f=\bar f j\), with \(j:X\hookrightarrow\bar X\) open and \(\bar f:\bar X\to Y\) proper (Theorem \bpout{A.2}). Since base change maps compose, it suffices to prove the lemma for \(j\) and for \(\bar f\) separately. For the open immersion \(j\), the base change map \(p_{\bar X}^*Rj_*M\to R(j\times\operatorname{id}_T)_*p_X^*M\) is an isomorphism by [Stacks, \href{https://stacks.math.columbia.edu/tag/0F1H}{0F1H}], applied to the cohomology sheaves of \(M\) and combined with the bounded-below spectral sequence. For \(\bar f\), proper base change, applied to \(Rj_*M\), shows that

\[
p_Y^*R\bar f_*Rj_*M\longrightarrow
R(\bar f\times\operatorname{id}_T)_*p_{\bar X}^*Rj_*M
\]

is an isomorphism. By Lemma \bpout{4.5} the composite of the two base change maps is the base change map for \(f\).

The step for \(j\) does not depend on Proposition \hyperref[bp64]{8.18} below, so there is no circularity. Indeed, the proof of [Stacks, \href{https://stacks.math.columbia.edu/tag/0F1H}{0F1H}] has two inputs. The first is base change under extension of the ground field [Stacks, \href{https://stacks.math.columbia.edu/tag/0F1C}{0F1C}], which by [Stacks, \href{https://stacks.math.columbia.edu/tag/0F09}{0F09}] follows from smooth base change (Theorem \hyperref[bp38]{8.9}), because after passing to perfect closures the larger field is a filtered union of smooth algebras over the smaller one. The second input, [Stacks, \href{https://stacks.math.columbia.edu/tag/0F1G}{0F1G}], is a Künneth formula with one proper factor, which comes from proper base change and the projection formula.
\end{proof}

% text: prose-08-19
\phantomsection\label{bp53}%
For \(f=j:\mathbf G_m\hookrightarrow\mathbf A^1\) and \(M=A\), the lemma says that the stalk of \(R(j\times\operatorname{id}_T)_*A\) at a point \((0,t)\) has the same cohomology as the stalk of \(Rj_*A\) at the origin, namely \(A\) in degree \(0\) and \(A(-1)\) in degree \(1\), whatever \(T\) is. The Künneth formula allows an arbitrary constructible second factor in place of the constant one, and we first define the map it concerns.

% object: definition-kunneth-map (blueprint/chap-Smooth-base-change/sec-The-Kunneth-formula/definition-kunneth-map/_section.md)
\phantomsection\label{bp54}%
\begin{bpstatement}{Definition~8.14}{ (Künneth map)}
Let \(f_i:X_i\to Y_i\), \(i=1,2\), be morphisms of separated \(k\)-schemes of finite type, put \(h=f_1\times f_2:X_1\times X_2\to Y_1\times Y_2\), and let \(M_i\in D^b_{\mathrm{ctf}}(X_i,A)\). The \emph{Künneth map}

\[
Rf_{1*}M_1\boxtimes Rf_{2*}M_2
\longrightarrow Rh_*(M_1\boxtimes M_2)
\tag{8.3}
\]

is the adjoint transpose, for \(h^*\dashv Rh_*\), of the morphism

\[
h^*(Rf_{1*}M_1\boxtimes Rf_{2*}M_2)
\simeq f_1^*Rf_{1*}M_1\boxtimes f_2^*Rf_{2*}M_2
\xrightarrow{\ \varepsilon_{f_1}\boxtimes\varepsilon_{f_2}\ }
M_1\boxtimes M_2,
\]

where \(\varepsilon_{f_i}\) is the counit of \(f_i^*\dashv Rf_{i*}\) and the first isomorphism comes from \(h^*p_i^*\simeq p_i^*f_i^*\).
\end{bpstatement}

% text: prose-08-20
\phantomsection\label{bp55}%
The case we prove first is the one in which one factor is an identity. For \(f:X\to Y\), \(T\) as in Lemma \hyperref[bp51]{8.13}, \(M\in D^b_{\mathrm{ctf}}(X,A)\) and \(Q\in D^b_{\mathrm{ctf}}(T,A)\), with \(f_1=f\) and \(f_2=\operatorname{id}_T\), put \(h=f\times\operatorname{id}_T\). The Künneth map is then

\[
\kappa_{M,Q}:Rf_*M\boxtimes Q
\longrightarrow Rh_*(M\boxtimes Q),
\tag{8.4}
\]

the adjoint transpose of \(h^*(Rf_*M\boxtimes Q)=p_X^*f^*Rf_*M\otimes^{\mathbf L} p_T^*Q
\xrightarrow{\varepsilon_f\otimes^{\mathbf L}1}M\boxtimes Q\). For constant \(Q\) it is the base change map of Lemma \hyperref[bp51]{8.13}, and every constructible \(Q\) is built from finite direct images of constant sheaves, which the next lemma treats.

% object: lemma-kunneth-finite-direct-image (blueprint/chap-Smooth-base-change/sec-The-Kunneth-formula/lemma-kunneth-finite-direct-image/_section.md)
\phantomsection\label{bp56}%
\begin{bpstatementit}{Lemma~8.15}{ (Finite direct images of constant sheaves)}
In the situation of (8.4), let \(Q\) be a finite direct sum of sheaves \(\pi_*A_{T'}\) with \(\pi:T'\to T\) finite. Then the Künneth map \(\kappa_{M,Q}\) is an isomorphism.
\end{bpstatementit}

% proof: proof/lemma-kunneth-finite-direct-image
\begin{proof}[{Proof}]\phantomsection\label{bp57}%
Both sides are additive in \(Q\), so we may assume \(Q=\pi_*A_{T'}\). Set

\[
v=\operatorname{id}_Y\times\pi,\qquad
w=\operatorname{id}_X\times\pi,\qquad
h'=f\times\operatorname{id}_{T'}.
\]

Direct image along a finite morphism is exact and commutes with pullback, and by the projection formula we have

\[
Rf_*M\boxtimes\pi_*A_{T'}
\simeq v_*p_{Y,T'}^*Rf_*M,
\qquad
M\boxtimes\pi_*A_{T'}\simeq w_*p_{X,T'}^*M.
\]

Since \(hw=vh'\), the comparison between these two expressions is \(v_*\) of the base change map of Lemma \hyperref[bp51]{8.13} for \(h'\), hence it is an isomorphism. Moreover, this comparison is \(\kappa_{M,Q}\), because under the adjunction \(h^*\dashv Rh_*\) and the adjunctions for the finite maps, both correspond to the counit \(\varepsilon_f\otimes^{\mathbf L}1\) that defines \(\kappa_{M,Q}\).
\end{proof}

% text: prose-08-21
\phantomsection\label{bp58}%
Over a field every constructible sheaf embeds into a sheaf of this form, and this allows an induction on the cohomological degree.

% object: lemma-kunneth-field-coefficients (blueprint/chap-Smooth-base-change/sec-The-Kunneth-formula/lemma-kunneth-field-coefficients/_section.md)
\phantomsection\label{bp59}%
\begin{bpstatementit}{Lemma~8.16}{ (Field coefficients)}
In the situation of (8.4), suppose that \(A=L\) is a finite field of characteristic \(\ell\). Then the Künneth map \(\kappa_{M,Q}\) is an isomorphism.
\end{bpstatementit}

% proof: proof/lemma-kunneth-field-coefficients
\begin{proof}[{Proof}]\phantomsection\label{bp60}%
We first treat constructible sheaves \(M\) and \(Q\), by the usual argument for comparing a \(\delta\)-functor with an exact functor, and then pass to complexes by truncation. We may compute with the small étale models of \(M\) and \(Q\), since Theorems \bpout{2.43} and \bpout{Theorem 2.38} identify the operations and the adjunction maps defining \(\kappa_{M,Q}\) on the two sites.

\emph{Sheaves.} Let \(M\) and \(Q\) be constructible sheaves of \(L\)-modules. By Stacks, [09Z7](https://stacks.math.columbia.edu/tag/09Z7), applied on \(T\), there is a short exact sequence of constructible \(L\)-sheaves

\[
0\longrightarrow Q\longrightarrow P\longrightarrow Q'\longrightarrow0,
\qquad
P=\bigoplus_a\pi_{a*}L_{T_a}^{\oplus n_a},
\quad \pi_a\text{ finite}.
\]

Indeed, every constructible sheaf on a qcqs scheme embeds into a finite sum of finite direct images of constant sheaves with finite coefficients, and over \(L\) these coefficients are finite sums of \(L\). For \(r\geq0\) put

\[
F^r(Q)=p_Y^*R^rf_*M\otimes_Lp_T^*Q,
\qquad
G^r(Q)=R^rh_*(p_X^*M\otimes_Lp_T^*Q).
\]

Since tensoring over \(L\) is exact, each \(F^r\) is exact in \(Q\), while the \(G^r\) form a long exact sequence. On cohomology sheaves the Künneth map induces natural maps \(a_Q^r:F^r(Q)\to G^r(Q)\), and \(a_P^r\) is an isomorphism for every \(r\) by Lemma \hyperref[bp56]{8.15}. It follows that each \(a_Q^r\) is injective. Indeed, in its naturality square for \(Q\hookrightarrow P\) the map \(F^r(Q)\to F^r(P)\) is injective and \(a_P^r\) is an isomorphism. Surjectivity is proved by induction on \(r\), simultaneously for all constructible \(Q\). For \(r=0\), the beginning of the long exact sequence of the \(G^r\) identifies \(G^0(Q)\) with the kernel of \(G^0(P)\to G^0(Q')\), and since \(F^0\) is exact, \(F^0(Q)\) is the kernel of \(F^0(P)\to F^0(Q')\). As \(a_P^0\) is an isomorphism and \(a_{Q'}^0\) is injective, \(a^0_Q\) maps the second kernel onto the first. For \(r>0\), the induction hypothesis for \(Q'\) and the surjection \(F^{r-1}(P)\to F^{r-1}(Q')\) imply that \(G^{r-1}(P)\to G^{r-1}(Q')\) is surjective. By the long exact sequence,

\[
G^r(Q)=\ker\bigl(G^r(P)\to G^r(Q')\bigr).
\]

Since \(F^r(Q)=\ker(F^r(P)\to F^r(Q'))\) by exactness of \(F^r\), the isomorphism \(a_P^r\) and the injectivity of \(a_{Q'}^r\) identify this kernel with \(F^r(Q)\). This completes the induction, so \(\kappa_{M,Q}\) is an isomorphism for sheaves.

\emph{Complexes.} The source and target of \(\kappa_{M,Q}\) are exact functors of \(M\) and of \(Q\), and \(\kappa_{M,Q}\) is natural in both. Using the finitely many truncation triangles of \(M\) and of \(Q\) and two-out-of-three, we can therefore extend the case of sheaves to \(M\in D^b_{\mathrm{ctf}}(X,L)\) and \(Q\in D^b_{\mathrm{ctf}}(T,L)\). This is possible because over the field \(L\) the truncations stay in \(D^b_{\mathrm{ctf}}\), since constructible cohomology sheaves have Tor amplitude zero, and all operations are bounded on the schemes under consideration.
\end{proof}

% text: prose-08-22
\phantomsection\label{bp61}%
Over a general finite ring the truncations of a complex in \(D^b_{\mathrm{ctf}}\) need not have finite Tor amplitude, so the argument by truncation does not apply directly. Instead we reduce to the residue fields and use derived Nakayama.

% object: lemma-kunneth-one-identity-factor (blueprint/chap-Smooth-base-change/sec-The-Kunneth-formula/lemma-kunneth-one-identity-factor/_section.md)
\phantomsection\label{bp62}%
\begin{bpstatementit}{Lemma~8.17}{ (Künneth formula with one identity factor)}
In the situation of (8.4), the Künneth map \(\kappa_{M,Q}\) is an isomorphism.
\end{bpstatementit}

% proof: proof/lemma-kunneth-one-identity-factor
\begin{proof}[{Proof}]\phantomsection\label{bp63}%
We show that the cone \(N\) of \(\kappa_{M,Q}\) vanishes modulo every maximal ideal of \(A\), and conclude by derived Nakayama on stalks. The source and target of \(\kappa_{M,Q}\) are in \(D^b_{\mathrm{ctf}}(Y\times T,A)\) by Proposition \bpout{7.1}, so every geometric stalk of \(N\) is a perfect \(A\)-complex. Now let \(\mathfrak m\subset A\) be a maximal ideal and put \(k=A/\mathfrak m\). The change-of-coefficients maps of Proposition \bpout{7.10} for pullback, tensor product and ordinary direct image respect the counit defining the Künneth map. Hence they identify the reduction \(\kappa_{M,Q}\otimes^{\mathbf L}_A k\) with \(\kappa_{M\otimes^{\mathbf L}_Ak,\,Q\otimes^{\mathbf L}_Ak}\), which is an isomorphism by Lemma \hyperref[bp59]{8.16}. This means that \(N\otimes^{\mathbf L}_A k=0\) for every maximal ideal. Since the Artinian ring \(A\) is a product of local rings, derived Nakayama for the perfect complex \(N_{\bar z}\) at a geometric point \(\bar z\) shows that \(N_{\bar z}=0\). Finally, since \(N\) lies in \(D^b_{\mathrm{ctf}}(Y\times T,A)\), it comes from a complex on the small étale site by Theorem \bpout{2.43}, where vanishing is detected on geometric stalks, so \(N=0\).
\end{proof}

% object: proposition-ordinary-external-kunneth (blueprint/chap-Smooth-base-change/sec-The-Kunneth-formula/proposition-ordinary-external-kunneth/_section.md)
\phantomsection\label{bp64}%
\begin{bpstatementit}{Proposition~8.18}{ (Ordinary external Künneth)}
Let \(f_i:X_i\to Y_i\), \(i=1,2\), be morphisms of separated \(k\)-schemes of finite type, and let \(M_i\in D^b_{\mathrm{ctf}}(X_i,A)\). Then the Künneth map (8.3)

\[
Rf_{1*}M_1\boxtimes Rf_{2*}M_2
\longrightarrow R(f_1\times f_2)_*(M_1\boxtimes M_2)
\]

is an isomorphism.
\end{bpstatementit}

% proof: proof/proposition-ordinary-external-kunneth
\begin{proof}[{Proof}]\phantomsection\label{bp65}%
We apply Lemma \hyperref[bp62]{8.17} once in each factor of the factorization

\[
X_1\times X_2\xrightarrow{\operatorname{id}_{X_1}\times f_2}
X_1\times Y_2\xrightarrow{f_1\times\operatorname{id}_{Y_2}}
Y_1\times Y_2
\]

of \(f_1\times f_2\). We apply it to \(f_1\) with second coefficient \(Rf_{2*}M_2\), which lies in \(D^b_{\mathrm{ctf}}(Y_2,A)\) by Proposition \bpout{7.1}. After interchanging the factors, we apply it to \(f_2\) with coefficient \(M_1\). Then the composite of \(R(f_1\times\operatorname{id}_{Y_2})_*\) of the second isomorphism with the first is an isomorphism from \(Rf_{1*}M_1\boxtimes Rf_{2*}M_2\) to \(R(f_1\times f_2)_*(M_1\boxtimes M_2)\). Since its adjoint transpose is \(\varepsilon_{f_1}\boxtimes\varepsilon_{f_2}\), it is the Künneth map (8.3).
\end{proof}



\clearpage\phantomsection\hypertarget{bp-outside}{}%
\section*{Outside this export}
\addcontentsline{toc}{section}{Outside this export}
References marked \ensuremath{{}^\ast} are to parts of the blueprint this document leaves out:
\begin{description}
\item[Chapter 2] The pro-étale site, in Ⅰ The six operations (\bpslug{chap-{}\allowbreak{}The-{}\allowbreak{}pro-{}\allowbreak{}etale-{}\allowbreak{}site})
\item[Corollary 2.37] Bounded-below classical complexes, in Ⅰ The six operations (\bpslug{corollary-{}\allowbreak{}bounded-{}\allowbreak{}below-{}\allowbreak{}classical-{}\allowbreak{}complexes})
\item[Theorem 2.38] Comparison of direct images, in Ⅰ The six operations (\bpslug{theorem-{}\allowbreak{}comparison-{}\allowbreak{}direct-{}\allowbreak{}images})
\item[Lemma 2.41] Torsors under classical groups, in Ⅰ The six operations (\bpslug{lemma-{}\allowbreak{}torsors-{}\allowbreak{}under-{}\allowbreak{}classical-{}\allowbreak{}groups})
\item[Theorem 2.43] Comparison for constructible sheaves, in Ⅰ The six operations (\bpslug{theorem-{}\allowbreak{}comparison-{}\allowbreak{}constructible-{}\allowbreak{}sheaves})
\item[Section 4.1] The base change map, in Ⅰ The six operations (\bpslug{sec-{}\allowbreak{}The-{}\allowbreak{}base-{}\allowbreak{}change-{}\allowbreak{}map})
\item[Lemma 4.5] Identity and pasting, in Ⅰ The six operations (\bpslug{lemma-{}\allowbreak{}identity-{}\allowbreak{}pasting})
\item[Lemma 4.7] Naturality, pasting and comparison, in Ⅰ The six operations (\bpslug{lemma-{}\allowbreak{}naturality-{}\allowbreak{}pasting-{}\allowbreak{}comparison})
\item[Theorem 4.20] Cohomology of curves, in Ⅰ The six operations (\bpslug{theorem-{}\allowbreak{}cohomology-{}\allowbreak{}curves})
\item[Lemma 4.27] Finite morphisms, in Ⅰ The six operations (\bpslug{lemma-{}\allowbreak{}finite-{}\allowbreak{}morphisms})
\item[Proposition 7.1] The six operations on $\Dctf$, in Ⅰ The six operations (\bpslug{proposition-{}\allowbreak{}six-{}\allowbreak{}operations-{}\allowbreak{}dctf})
\item[Proposition 7.10] The six operations and change of coefficients, in Ⅰ The six operations (\bpslug{proposition-{}\allowbreak{}six-{}\allowbreak{}operations-{}\allowbreak{}change-{}\allowbreak{}coefficients})
\item[Proposition 7.12] Isomorphisms reduce to finite coefficients, in Ⅰ The six operations (\bpslug{proposition-{}\allowbreak{}isomorphisms-{}\allowbreak{}reduce-{}\allowbreak{}finite-{}\allowbreak{}coefficients})
\item[Chapter 9] Smooth purity, in Ⅱ Duality and the trace formula (\bpslug{chap-{}\allowbreak{}Smooth-{}\allowbreak{}purity})
\item[Theorem A.2] Nagata's compactification theorem, in Appendices (\bpslug{theorem-{}\allowbreak{}nagata-{}\allowbreak{}s-{}\allowbreak{}compactification-{}\allowbreak{}theorem})
\item[Appendix D] Cohomology of torsion sheaves on curves, in Appendices (\bpslug{chap-{}\allowbreak{}Cohomology-{}\allowbreak{}of-{}\allowbreak{}torsion-{}\allowbreak{}sheaves-{}\allowbreak{}on-{}\allowbreak{}curves})
\end{description}
\clearpage\phantomsection%
\section*{References}
\addcontentsline{toc}{section}{References}
\bpbibitem{BBD82}{bp67}{A. Beilinson, J. Bernstein, and P. Deligne. \emph{Faisceaux pervers}, volume 100 of \emph{Astérisque}. 1982.}

\bpbibitem{BS15}{bp68}{B. Bhatt and P. Scholze. The pro-étale topology for schemes. \emph{Astérisque}, 369:99–201, 2015. \href{https://people.mpim-bonn.mpg.de/scholze/proetale.pdf}{Author's manuscript}.}

\bpbibitem{Con07}{bp69}{B. Conrad. Deligne's notes on Nagata compactifications. \emph{Journal of the Ramanujan Mathematical Society}, 22:205–257, 2007. \href{https://math.stanford.edu/\~conrad/papers/nagatafinal.pdf}{Author's manuscript}.}

\bpbibitem{Del74}{bp70}{P. Deligne. La conjecture de Weil. I. \emph{Publ. Math. IHÉS}, 43:273–307, 1974.}

\bpbibitem{Del80}{bp71}{P. Deligne. La conjecture de Weil. II. \emph{Publ. Math. IHÉS}, 52:137–252, 1980.}

\bpbibitem{Den22}{bp72}{C. Deninger. There is no “Weil-”cohomology theory with real coefficients for arithmetic curves. \emph{Ann. Sc. Norm. Super. Pisa Cl. Sci.}, 2022. To appear. arXiv:\href{https://arxiv.org/abs/2204.02714}{2204.02714}.}

\bpbibitem{Deu41}{bp73}{M. Deuring. Die Typen der Multiplikatorenringe elliptischer Funktionenkörper. \emph{Abh. Math. Sem. Univ. Hamburg}, 14:197–272, 1941.}

\bpbibitem{Eke90}{bp74}{T. Ekedahl. On the adic formalism. In \emph{The Grothendieck Festschrift, Vol. II}, volume 87 of \emph{Progress in Mathematics}, pages 197–218. Birkhäuser, 1990.}

\bpbibitem{Has36}{bp75}{H. Hasse. Zur Theorie der abstrakten elliptischen Funktionenkörper. I, II, III. \emph{J. Reine Angew. Math.}, 175:55–62, 69–88, 193–208, 1936.}

\bpbibitem{Hum72}{bp76}{J. E. Humphreys. \emph{Introduction to Lie Algebras and Representation Theory}, volume 9 of \emph{Graduate Texts in Mathematics}. Springer, 1972.}

\bpbibitem{Ill94}{bp77}{L. Illusie. Autour du théorème de monodromie locale. \emph{Astérisque}, 223:9–57, 1994. §4. \href{https://www.numdam.org/item/AST_1994__223__9_0/}{Original article}.}

\bpbibitem{Kat88}{bp78}{N. M. Katz. Travaux de Laumon. \emph{Astérisque}, 161–162:105–132, 1988. Séminaire Bourbaki, exp. 691.}

\bpbibitem{KW01}{bp79}{R. Kiehl and R. Weissauer. \emph{Weil conjectures, perverse sheaves and \(\ell\)-adic Fourier transform}, volume 42 of \emph{Ergebnisse der Mathematik}. Springer, 2001.}

\bpbibitem{Lau87}{bp80}{G. Laumon. Transformation de Fourier, constantes d'équations fonctionnelles et conjecture de Weil. \emph{Publ. Math. IHÉS}, 65:131–210, 1987.}

\bpbibitem{Mil86a}{bp81}{J. S. Milne. \href{https://www.jmilne.org/math/xnotes/AVs.pdf}{Abelian Varieties}. In \emph{Arithmetic Geometry}. 1986. Corrected online version.}

\bpbibitem{Mil86b}{bp82}{J. S. Milne. \href{https://www.jmilne.org/math/xnotes/JVs.pdf}{Jacobian Varieties}. In \emph{Arithmetic Geometry}. 1986. Corrected online version.}

\bpbibitem{Ran39}{bp83}{R. A. Rankin. Contributions to the theory of Ramanujan's function \(\tau(n)\) and similar arithmetical functions. II. \emph{Proc. Cambridge Philos. Soc.}, 35:357–372, 1939.}

\bpbibitem{SGA4}{bp66}{M. Artin, A. Grothendieck, and J.-L. Verdier, editors. \emph{Théorie des topos et cohomologie étale des schémas}, volumes 269, 270, 305 of \emph{Lecture Notes in Mathematics}. Springer, 1972–1973. SGA 4.}

\bpbibitem{SGA4½}{bp84}{P. Deligne. \emph{Cohomologie étale}, volume 569 of \emph{Lecture Notes in Mathematics}. Springer, 1977. SGA 4½; \emph{Théorèmes de finitude}, §3, and \emph{Rapport sur la formule des traces}, §3, and \emph{Sommes trigonométriques}, §1.}

\bpbibitem{SGA5}{bp85}{A. Grothendieck, editor. \emph{Cohomologie \(\ell\)-adique et fonctions \(L\)}, volume 589 of \emph{Lecture Notes in Mathematics}. Springer, 1977. SGA 5.}

\bpbibitem{SGA7II}{bp86}{P. Deligne and N. Katz, editors. \emph{Groupes de monodromie en géométrie algébrique}, volume 340 of \emph{Lecture Notes in Mathematics}. Springer, 1973. SGA 7 II; Exposé XIII, (2.1.7).}

\bpbibitem{Stacks}{bp87}{The Stacks Project Authors. \href{https://stacks.math.columbia.edu}{The Stacks Project}. Individual results are cited by their stable tags in the text.}

\bpbibitem{Var05}{bp88}{Y. Varshavsky. Lefschetz–Verdier trace formula and a generalization of a theorem of Fujiwara. 2005. Version 2. arXiv:\href{https://arxiv.org/abs/math/0505564}{math/0505564}. \href{https://arxiv.org/abs/math/0505564v2}{Author's preprint}.}

\bpbibitem{Wei48a}{bp89}{A. Weil. \emph{Sur les courbes algébriques et les variétés qui s'en déduisent}. Hermann, Paris, 1948.}

\bpbibitem{Wei48b}{bp90}{A. Weil. \emph{Variétés abéliennes et courbes algébriques}. Hermann, Paris, 1948.}

\bpbibitem{Wei49}{bp91}{A. Weil. Numbers of solutions of equations in finite fields. \emph{Bull. Amer. Math. Soc.}, 55:497–508, 1949.}

\end{document}
